% HUSTthesis.tex @ https://github.com/zfengg/HUSTtex/tree/master/HUSTthesis
%
% 	A simple tex template for the undergraduate thesis at HUST.
%
% author: Zhou Feng @ 2019
%
% requirements:
%	TeX environments: TeXlive/MacTeX or MiKTeX
%	compiler: XeLaTeX
%
% ---------------------------------------------------------------------------- %
%                                   preamble                                   %
% ---------------------------------------------------------------------------- %
\documentclass[a4paper]{article}
% 10.5pt = 5 hao

% ---------------------------------- layout ---------------------------------- %
\usepackage[a4paper,top=2.5cm,bottom=2.5cm,left=3cm,right=3cm,% margins
			headheight=1.5cm,headsep=1.5em,
			footskip=2em,
			]{geometry}

% ------------------------------- math symbols ------------------------------- %
% 载入常用的数学包, 符号包
\usepackage{amsmath}
\usepackage{amsfonts}
\usepackage{amssymb}
\usepackage{mathrsfs}
\usepackage{blindtext}

%----------------------------------------------------------------%
%% linespace 行间距，段间距等等
\usepackage{setspace}
% \usepackage{indentfirst} % then the first line of each title should start with a indent.
% 定义标题和段落样式
% 定义1.5倍行距
\renewcommand{\baselinestretch}{1.62}
\setlength{\baselineskip}{12pt}   % set the fixed value of the lineskip
\setlength\parskip{\baselineskip} % set the space between the paragraphs, set the variable \parskip \baselineskip
% parindent
\setlength{\parindent}{0pt}

%----------------------------------------------------------------%
% fonts (style, color, size). 字体的大小，颜色，以及定义常用的字号
\usepackage{ctex}		 	% If you are lazy, the CTEX suit is enough.
% Chinese font
\usepackage{xeCJK}		 	% For the Chinese through XeLaTex
\setCJKmainfont{SimSun} 	% set the mainfont of Chinese as songti. (serif) for
\setCJKsansfont{SimSun} 	% sans serif font for \textsf
\setCJKmonofont{SimSun} 	% monospace font for \texttt
% \punctstyle{kaiming}  	% Remove the space used by symbols like comma. Special for the mainland students like us HUSTers.
\setCJKfamilyfont{song}{SimSun}                       % 宋体 song
\newcommand{\song}{\CJKfamily{song}}
\setCJKfamilyfont{kai}{KaiTi}                         % 楷体  kai
\newcommand{\kai}{\CJKfamily{kai}}
\setCJKfamilyfont{hwzs}{STZhongsong}                  % 华文中宋  hwzs
\newcommand{\hwzs}{\CJKfamily{hwzs}}
\setCJKfamilyfont{hei}{SimHei}                        % 黑体  hei
\newcommand{\hei}{\CJKfamily{hei}}
% English font
\usepackage{fontspec}% Then you can use the fonts installed at your device.
\setmainfont{Times New Roman}
\setsansfont{Times New Roman}
\setmonofont{Times New Roman}
%\setsansfont{[foo.ttf]} % for the fonts at this default path.
% font Color 利用definecolor自己可以定义颜色
\usepackage{xcolor}
\definecolor{MSBlue}{rgb}{.204,.353,.541}
\definecolor{MSLightBlue}{rgb}{.31,.506,.741}
% font Size (I use pinyin represents the corresponding size in Microsorft Word)
% \newcommand{\chuhao}{\fontsize{42pt}{\baselineskip}\selectfont}
\newcommand{\xiaochuhao}{\fontsize{36pt}{\baselineskip}\selectfont}
% \newcommand{\yihao}{\fontsize{28pt}{\baselineskip}\selectfont}
\newcommand{\erhao}{\fontsize{21pt}{\baselineskip}\selectfont}
\newcommand{\xiaoerhao}{\fontsize{18pt}{\baselineskip}\selectfont}
\newcommand{\sanhao}{\fontsize{15.75pt}{\baselineskip}\selectfont}
\newcommand{\sihao}{\fontsize{14pt}{18pt}\selectfont}
\newcommand{\xiaosihao}{\fontsize{12pt}{18pt}\selectfont}
\newcommand{\wuhao}{\fontsize{10.5pt}{18pt}\selectfont}
% \newcommand{\xiaowuhao}{\fontsize{9pt}{\baselineskip}\selectfont}
% \newcommand{\liuhao}{\fontsize{7.875pt}{\baselineskip}\selectfont}
% \newcommand{\qihao}{\fontsize{5.25pt}{\baselineskip}\selectfont}

% ---------------------------------------------------------------------------- %
%% header and footer 页眉，页脚
\usepackage{fancyhdr} % for header and footer
% 设置页眉样式
\newcommand{\headstyle}{
 \fancyhead[C]{ \hwzs\wuhao 华\hspace{0.5em}中\hspace{0.5em}科\hspace{0.5em}技\hspace{0.5em}大\hspace{0.5em}学\hspace{0.5em}毕\hspace{0.5em}业\hspace{0.5em}设\hspace{0.5em}计\hspace{0.5em}(论\hspace{0.5em}文)}
}
% 设置页脚样式
\newcommand{\footstyle}{\fancyfoot[C]{\wuhao\thepage}
 \fancyfoot[L]{\rule[5pt]{6.7cm}{0.4pt}}
 \fancyfoot[R]{\rule[5pt]{6.7cm}{0.4pt}}
}
\pagestyle{fancy}
\fancyhf{} % 清空原有样式
\headstyle
\footstyle
% 定义一种新的格式叫做main
\fancypagestyle{main}{%
 \fancyhf{} % 清空原有样式
 \headstyle
 \footstyle
}
\renewcommand{\headrulewidth}{0.4pt}
% \renewcommand{\footrulewidth}{0.4pt}
% \renewcommand{\headrule}{\rule{\textwidth}{0.4pt}}

% ---------------------------------------------------------------------------- %
% set the styles of sections at all levels
% 设置各个标题样式
% 不需要使用part和chapter层级
\usepackage{titlesec}
\usepackage{titletoc}
\titleformat{\section}{\centering\hei\bfseries\xiaoerhao}{\thesection}{1em}{} % 在section标题编号后面加个点
% \titlelabel{\thetitle.\quad} % add a dot after the counter for all levels of sections
% \titleformat*{\section}{\wuhao\bfseries} % 设置标签的形式，5号加粗
\titleformat*{\subsection}{\raggedright\hei\bfseries\sihao}
\titleformat*{\subsubsection}{\raggedright\hei\bfseries\xiaosihao}
\titleformat{\paragraph}[hang]{\raggedright\hei\bfseries\xiaosihao}{\theparagraph}{1em}{}[]

% manual
% \titleformat{command}[shape]{format}{label}{sep}{before-code}[after-code]
% \titlespacing{command}{left}{before-sep}{after-sep}
% 设置新的层级subsubsubsection
\setcounter{tocdepth}{4}
\setcounter{secnumdepth}{4}

\newcommand{\sectionbreak}{\clearpage} % 小节从新的一页开始
% 根据学校要求设置新的section, subsection, subsection,  paragraph

% set the content of section and so on
\newcommand\seccontent{
	\song
	\xiaosihao % 默认五号字体, 行间距为1.5*\baselineskip
    \setlength{\parindent}{2em} % 首段缩进两个M字符
    \setlength{\parskip}{0pt}
    }
\newcommand\tabcontent{
	\song
	\wuhao % 默认五号字体, 行间距为1.5*\baselineskip
	\setlength{\parindent}{2em} % 首段缩进两个M字符
	\setlength{\parskip}{0pt}
}


% ---------------------------------------------------------------------------- %
% for the style of theorems, definitions, proofs and remarks 定义数学里面一些常用的环境
\usepackage{amsthm}
\newtheorem{thm}{\textbf{定理}}[section]
% The section in [] can be replaced by chapter or subsection
\theoremstyle{definition} 
\theoremstyle{plain}
\theoremstyle{remark}

% ---------------------------------------------------------------------------- %
% for the caption and reference 图表及公式的编号规范
\usepackage{float} 		 		  	% table figure positioning
\usepackage{caption}
\captionsetup[figure]{labelformat=default, labelsep=quad,name={图}}
\captionsetup[table]{labelformat=default,labelsep=quad,name={表}}
% 设置图表标题的计数方式
\renewcommand{\thefigure}{\ifnum \thesection>0 \thesection-\fi \arabic{figure}} % set caption label style to 2-1
\renewcommand{\thetable}{\ifnum \thesection>0 \thesection-\fi \arabic{table}} % set caption label style to 2-1
\DeclareCaptionFont{mylabelfont}{\hei\xiaosihao}
\captionsetup[figure]{font=mylabelfont}
\captionsetup[table]{font=mylabelfont}

% 设置图表的autoref的格式
\newcommand{\reffig}[1]{图 \ref{#1}}
\newcommand{\reftab}[1]{表 \ref{#1}}
% 公式的编号格式
\renewcommand\theequation{\arabic{section}-\arabic{equation}}

\usepackage{graphicx} % To include graphixs 添加图所需的包
\usepackage{booktabs} % To create three line table including the commands toprule, bottomrule, and midrule
% \usepackage{colortbl} %
% 使用tabularx库并定义新的左右中格式
\usepackage{tabularx}
\usepackage{makecell}
\newcolumntype{L}{X}
\newcolumntype{C}{>{\centering \arraybackslash}X}
\newcolumntype{R}{>{\raggedright \arraybackslash}X}

% ---------------------------------------------------------------------------- %
% set the style of counters 设置计数器
% 设置重新计数的位置
\makeatletter
\@addtoreset{footnote}{page}
\@addtoreset{figure}{section}
\@addtoreset{table}{section}
\@addtoreset{equation}{section}
\makeatother

% ---------------------------------------------------------------------------- %
% tableofcontents, listoftables and listoffigures 目录
%\renewcommand\listfigurename{插图列表}
%\renewcommand\listtablename{表格列表}
%\titlecontents{标题名}[左间距]{标题格式}{标题标志}{无序号标题}{指引线与页码}[下间距]
%\dottedcontents{section}[2.55em]{\song \xiaosihao \bfseries}{2.5em}{1em}
\usepackage{tocloft}
\renewcommand{\contentsname}{\centerline{ \hei\bfseries\xiaoerhao 目\hspace{2em}录}}
\titlecontents{section}[3em]{\song\xiaosihao\bfseries}{\contentslabel{3em}}{\hspace*{-3em}}{\normalfont\titlerule*[8pt]{.}\contentspage}
\titlecontents{subsection}[3em]{\song\xiaosihao}{\contentslabel{3em}}{\hspace*{-3em}}{\titlerule*[8pt]{.}\contentspage}
\titlecontents{subsubsection}[4em]{\song\xiaosihao}{\contentslabel{4em}}{\hspace*{-4em}}{\titlerule*[8pt]{.}\contentspage}
\titlecontents{paragraph}[5em]{\song\xiaosihao}{\contentslabel{5em}}{\hspace*{-5em}}{\titlerule*[8pt]{.}\contentspage}

% ---------------------------------------------------------------------------- %
% reference and citation 参考文献
\usepackage{natbib}
\renewcommand{\refname}{\centering\hei\xiaoerhao 参考文献}
\bibsep=0pt % 用来设置每个\bibitem之间的间距
% \renewcommand{\bibname}{HUSTthesis} % .bib name
% \newcommand{\upcite}[1]{\textsuperscript{\textsuperscript{\cite{#1}}}} % show citation label in the upperscript

% ---------------------------------------------------------------------------- %
% 设置声明页
% 使用特殊符号
\usepackage{amssymb}
\usepackage{wasysym}
% 制作tatement中的符号
\def\HUSTcheckedbox{$\Square\!\!\!\!\checkmark$}
\def\HUSTbox{$\Square$}
% 设置命令
\newcommand{\statement}[2]{
	\def\yearofconfidentiality{\hspace{2em}}

    \def\HUSTconfidential{\HUSTbox}
	\def\HUSTnotconfidential{\HUSTcheckedbox}
	\ifnum #1<2
		\def\HUSTconfidential{\HUSTcheckedbox}
		\def\HUSTnotconfidential{\HUSTbox}
		\def\yearofconfidentiality{#2}
	\fi
	\clearpage
	\thispagestyle{main}
	\vspace*{1em}
	\seccontent
	\begin{center}
		\heiti \xiaoerhao \bfseries
		学位论文原创性声明
	\end{center}


    本人郑重声明：所呈交的论文是本人在导师的指导下独立进行研究所取得的研究成果。除了文中特别加以标注引用的内容外，本论文不包括任何其他个人或集体已经发表或撰写的成果作品。本人完全意识到本声明的法律后果由本人承担。

	\begin{flushright}
		作者签名：\hspace{6em} 年 \hspace{2em} 月 \hspace{2em} 日
	\end{flushright}
	\vspace{4em}

	\begin{center}
		\heiti \zihao{-2} \bfseries
		学位论文版权使用授权书
	\end{center}


    本学位论文作者完全了解学校有关保障、使用学位论文的规定，同意学校保留并向有关学位论文管理部门或机构送交论文的复印件和电子版，允许论文被查阅和借阅。本人授权省级优秀学士论文评选机构将本学位论文的全部或部分内容编入有关数据进行检索，可以采用影印、缩印或扫描等复制手段保存和汇编本学位论文。

    本学位论文属于 1、保密 \HUSTconfidential ，在 \yearofconfidentiality 年解密后适用本授权书。

		    \hspace{7em} 2、不保密 \HUSTnotconfidential \hspace{1em}
			。

			\hspace{7em} (请在以上相应方框内打``$\checkmark$'')



	\begin{flushright}
	作者签名：\hspace{6em} 年 \hspace{2em} 月 \hspace{2em} 日

    导师签名：\hspace{6em} 年 \hspace{2em} 月 \hspace{2em} 日
    \end{flushright}
	\clearpage
}
\newcommand{\makestatement}[2]{\statement{#1}{#2}}

% ---------------------------------------------------------------------------- %
% 定义中英文摘要和致谢环境
%
% 中文摘要环境
\newenvironment{cnabstract}[1]{
	\def \cnkeyword {#1}
	\clearpage
	\phantomsection
	\addcontentsline{toc}{section}{摘要}
	\vspace*{-20pt}
	\begin{center}
		\heiti \bfseries \xiaoerhao 摘 \hspace{2em} 要
	\end{center}
	\seccontent
}{
	\vspace{1em}
	\par\noindent {\hei\sihao \bfseries 关键词：} {\song\xiaosihao\cnkeyword}

}

% 英文摘要环境
\newenvironment{enabstract}[1]{
	\def \enkeyword {#1}
	\clearpage
	\phantomsection
	\addcontentsline{toc}{section}{Abstract}
	\vspace*{-20pt}
	\begin{center}
		\bfseries \xiaoerhao Abstract
	\end{center}
	\seccontent
}{
	\vspace{1em}
	\par\noindent {\sihao\bfseries Key Words: }\ {\song\xiaosihao\enkeyword}
	\clearpage
}

% 定义致谢环境
\newenvironment{thankpage}{
	\clearpage%
	\phantomsection%
	\addcontentsline{toc}{section}{致谢}%
	\section*{致\hspace{2em}谢}%
}{
	\clearpage
}


% ---------------------------------------------------------------------------- %
%	---	定义列表项，列举的样式
\usepackage{enumitem}
\setlist{noitemsep}

% ---------------------------------------------------------------------------- %
% \usepackage{makeindex} For the index 索引
\usepackage{listings} % For the code. 代码

% ---------------------------------------------------------------------------- %
% 设置脚注

% ---------------------------------------------------------------------------- %
%% For the hyperlink and bookmark 超链接及书签，(这样生成的pdf中的引用直接点击链接即可到达目的地)
\usepackage[bookmarks=true,colorlinks,linkcolor=black,citecolor=black,urlcolor=purple]{hyperref}% 设置超链接并修改风格


% ---------------------------------------------------------------------------- %
%% For the appendix, 附录
% 设置附录
\usepackage{appendix}
\renewcommand{\appendixname}{附录}

% ---------------------------------------------------------------------------- %
% for the titlepage 标题页，此处可以省略，建议直接使用官方给出的标题页即可
\usepackage{titling}
% 重置命令 maketitle
\renewcommand{\maketitle}{
	\def\HUSTtitlelength{12em}
 	\begin{titlepage}
		\begin{center}
			\vspace*{0em}
			% \includegraphics[height=1.61cm]{HUSTlogo.eps}\\
			\includegraphics[height=1.61cm]{HUSTGreen.eps}\\
%
			\vspace*{4em}
%
			{\xiaochuhao \hwzs \bfseries 本科生毕业设计[论文]}\\
%
			\vspace*{6em}
			{\erhao \hei \bfseries \thetitle}

			\vspace*{6em}
			{\sanhao \hwzs
				\renewcommand\arraystretch{2.7}
				\begin{tabular}{lc}
					\makebox[4em][s]{院 \hfill 系} &
					\underline{\makebox[\HUSTtitlelength]{\school}} \\
					\makebox[4em][s]{专业班级} &
					\underline{\makebox[\HUSTtitlelength]{\classnum}} \\
					\makebox[4em][s]{姓 \hfill 名} &
					\underline{\makebox[\HUSTtitlelength]{\theauthor}} \\
					\makebox[4em][s]{学 \hfill 号} &
					\underline{\makebox[\HUSTtitlelength]{\stunum}} \\
					\makebox[4em][s]{指导教师} &
					\underline{\makebox[\HUSTtitlelength]{\instructor}} \\
			  \end{tabular}
		    }

			\vspace{4em}
			{\sanhao \hwzs \thedate}

		\end{center}
	\end{titlepage}
}

% ------------------------------------ 标题页 ----------------------------------- %
\title{连分数动力系统的混合性} % 论文题目
\def\school{} % 院系
\def\classnum{} % 专业班级
\author{王子雄} % 姓名
\def\stunum{}	% 学号
\def\instructor{} % 指导老师
\date{2023年5月24日} % 日期

% -------------------------------- quickinput -------------------------------- %
% 利用\newcommand{cmd}{def} 设置一些常用的代码，提高效率，这里可以自行删除，下面是我敲翻译时候打的一些command。

% ---------------------------------------------------------------------------- %
%                                   document                                   %
% ---------------------------------------------------------------------------- %
\begin{document}

\maketitle

% \thispagestyle{empty}% 标题页不参与编号
% ------------------------------------ 声明页 ----------------------------------- %
\setcounter{page}{1}
\renewcommand{\thepage}{\Roman{page}}
\makestatement{2}{2019} % 生成声明页,第一个参数选择1则加上保密的时间，此处为2019，如果选择2则选择不保密。

% ----------------------------------- 中文摘要 ----------------------------------- %
\setcounter{page}{1}
\renewcommand{\thepage}{\Roman{page}}
\begin{cnabstract}{连分数；Kuzmin定理；Gauss测度；混合性；遍历性}
    \hspace{2em}
    本文主要讨论了连分数动力系统的强混合性与遍历性,并由此得到了部分商序列的渐进独立性, $\psi$-混合性以及一些极限性质.\par
    第一章简单介绍了连分数的研究历史和现状.\par
    第二章介绍了连分数的定义和基本性质,得到了连分数收敛的判定条件,从而说明了正则连分数的收敛性.之后介绍了连分数度量理论的基本内容,并研究了Gauss变换和Gauss测度,说明了Gauss 变换保持 Gauss 测度以及 Gauss 测度与 Lebesgue 测度的等价性,连分数部分商序列的严平稳、可逆和不独立性.\par
    第三章介绍了保测变换和半代数、代数、单调类,并给出了保测变换的验证方法.之后,通过Kuzmin定理给出了Gauss问题的一个解以及证明了连分数动力系统的强混合性与部分商序列的$\psi$-混合性、渐进独立性.\par
    第四章介绍了遍历系统和遍历论的基本定理——Birkhoff遍历定理.之后,证明了强混合性蕴含弱混合性蕴含遍历性,并由此证明了连分数动力系统的遍历性.并且运用Birkhoff遍历定理证明了部分商序列的三个性质.

\end{cnabstract}

% ----------------------------------- 英文摘要 ----------------------------------- %
\begin{enabstract}{Continued Fractions; Kuzmin's Theorem; Gauss Measure; Mixing Property; Ergodicity}

This paper mainly discusses the strong mixing property and ergodicity of the continued fraction dynamical system, and thus obtains the asymptotic independence, $\psi$-mixing, and some limit properties of the partial quotient sequence.

Chapter 1 provides a brief introduction to the research history and current state of continued fractions.

Chapter 2 expounds on the definition and fundamental properties of continued fractions, establishes a convergence test of continued fractions, and thereby demonstrates the convergence of regular continued fractions. Subsequently, the basic principles of the metric theory of continued fractions, the Gauss transformation and Gauss measure are introduced. It is demonstrated that the Gauss transformation preserves the Gauss measure and that the Gauss measure is equivalent to the Lebesgue measure. Additionally, the strict stationarity, reversibility, and non-independence of the partial quotient sequence of continued fractions are proved.

Chapter 3 introduces measure-preserving transformations, semi-algebra, algebra, and monotonic classes, then provides a method for verifying measure-preserving transformations. Subsequently, using Kuzmin's theorem, a solution to the Gauss problem is given, and the strong mixing of the continued fraction dynamical system, as well as the $\psi-$mixing and asymptotic independence of the partial quotient sequence, are proven.

Chapter 4 introduces ergodic systems and the fundamental theorem of ergodic theory - Birkhoff's ergodic theorem. After that, by proving that strong mixing implies ergodicity and ergodicity implies weak mixing, it proves the ergodicity of the continued fraction dynamical system. Additionally, by applying Birkhoff's ergodic theorem three properties of the partial quotient sequence are proved.

\end{enabstract}


\vspace*{-1em}
\tableofcontents
\thispagestyle{main}


\clearpage
\setcounter{page}{1}
\renewcommand{\thepage}{\arabic{page}}

% ----------------------------------- 主体内容 ----------------------------------- %
\seccontent
\section{绪论}
\hspace{12em}

将形如
\begin{equation}\label{definition}
      a_0+\frac{1}{a_1+\dfrac{1}{a_2+\cdots+\dfrac{1}{a_n+\cdots}}}
\end{equation}
的式子称为连分数.关于连分数的研究历史悠久,在古希腊数学家Euclid计算两个正整数$a,b$的最大公因数的辗转相除法中,实际上已经得出了有理数$\dfrac{a}{b}$的连分数表示.在另外两个古老的问题:平方根估计和丢番图方程$ax+by=c$求解中,早期的数学家们也在某种意义上使用了连分数.在这之后,意大利数学家Fibonacci首次尝试对连分数进行一般定义,而Bombelli被认为是真正使用了现代连分数理论的第一位数学家.到了18世纪, Euler, Lambert和Lagrange等数学家为连分数理论做出了重大的贡献.在19世纪初, Gauss提出的一个问题让连分数从纯粹的数论研究领域变成结合现代分析的研究领域.这个问题出现在Gauss写给Laplace的信中, Gauss想要估计
$$m_n(x)-\frac{\ln(1+x)}{\ln{2}}$$
Gauss在信中声称他已经证明了$\lim\limits_{n\to\infty}m_n(x)=\dfrac{\ln(1+x)}{\ln{2}},\quad(0\leqslant x\leqslant 1)$,其中
$m_n(x)=\lambda({u\in[0,1)\mid\tau^n(u)\leqslant x})$, $\lambda$为Lebesgue测度, $\tau$为Gauss变换.\par
一个世纪后, Kuzmin给出了上面问题的一个估计\textsuperscript{\cite{Kuzmin}}:
$$m_n\left(x\right)=\frac{\ln{\left(1+x\right)}}{\ln{2}}+r_n\left(x\right),\quad r_n\left(x\right)=O\left(q^{\sqrt n}\right).$$\par
Levy用概率语言给出了更精细的估计\textsuperscript{\cite{Iosifescu,Levy}}:
$$r_n(x)=O(q^n).$$\par
从此,对这个问题以及相关问题的研究得到高速发展,详情请见参考文献\cite{Liu_Peng,The_long_history,Iosifescu,Gauss-Kuzmin_type,Zhang_Xian,Random_variables,Dimensional_theory,Zhang_Mengjie,Xie_Shenghan,Geng_Xianjin},并在纯数学和自然科学等多个领域中得到了比较多的成果.目前对连分数的研究主要涉及到动力系统,遍历论,概率论,随机过程,分形几何等知识.
在本文中,我们将给出连分数收敛的判定定理;解释连分数部分商序列严平稳,可逆等性质;介绍混合性与遍历性的基本知识;并证明连分数动力系统的强混合性;由强混合性与遍历性的相关关系证明连分数动力系统的遍历性,并进行简单运用.\par
本论文概要如下:
\begin{enumerate}
      \item 绪论,主要介绍连分数的研究历史和现状,并介绍文章主要结论.
      \item 对连分数的定义、收敛性、度量理论以及Gauss变换和Gauss测度以及部分商序列的一些性质进行介绍.
      \item 对动力系统混合性的介绍,证明Kuzmin定理,并由此解决Gauss问题及证明连分数动力系统的强混合性,连分数部分商序列的$\psi$-混合性.
      \item 对动力系统遍历性的介绍,由连分数动力系统的强混合性证明连分数动力系统的遍历性,并给出由遍历性得到的关于连分数部分商序列的结论.
\end{enumerate}

\sectionbreak
% ---------------------------------------------------------------------------- %
\section{连分数简介}
\subsection{连分数的定义与收敛性}

\textbf{定义 2.1.1.  }
将连分数展式(\ref{definition})简记为
$$[a_0;a_1,a_2,\cdots,a_n,\cdots],$$
并称$a_1,a_2,\cdots$为它的部分商.\par
\par
其中,部分商$a_1,a_2,\cdots$可以取在实数域$\mathbb{R}$,复数域$\mathbb{C}$等数域,甚至一些抽象的函数空间中. $a_0$为整数, $a_1,a_2,\cdots$均为正整数的一类连分数有重要的研究意义,我们称其为正规连分数.在本文中,我们主要研究正规连分数.不过,为得到连分数收敛性的一个判定定理,我们先仅仅限制部分商为正的.\par
一个连分数称为有限的,若其只有有限多个部分商;称为无限的,若其有无穷多个部分商.\par
\textbf{定义 2.1.2.  }
连分数的$n$阶收敛因子为
$$\frac{p_n}{q_n}=[a_0;a_1,a_2,\cdots,a_n].$$
\par
考虑收敛因子与部分商的关系,由于$n$阶收敛因子是前$n$个部分商及$a_0$进行有理运算的结果,所以它是前$n$个部分商及$a_0$的有理函数,即
$$\frac{p_n}{q_n}=\frac{P(a_0,a_1,\cdots ,a_n)}{Q(a_0,a_1,\cdots ,a_n)}.$$\par
上式也给出了有限连分数的值.而对无限连分数,我们需要考虑它的收敛性,为此,对收敛因子的研究是重要的.\par
\textbf{命题 2.1.1.  }\textsuperscript{\cite{Khinchin}}
$\forall k\geqslant2$都有
    \begin{equation}\label{rfc}
        \begin{aligned}
        p_k=a_kp_{k-1}+p_{k-2},\\
        q_k=a_kq_{k-1}+q_{k-2}.
        \end{aligned}
    \end{equation}
若我们约定$p_{-1}=1,q_{-1}=0$则\textbf{命题 2.1.1}还将对$k=1$成立.
\par
\textbf{证明:  }
我们将采用归纳的方式证明, $k=2$时可以直接验证,下设该命题对$k<n$都成立.\par
记$[a_1;a_2,\cdots,a_n]$的$r$阶收敛因子为$[a_1;a_2,\ldots,a_r]=\dfrac{p_r^{'}}{q_r^{'}},$则
$$[a_0;a_1,\cdots,a_n]=a_0+\dfrac{q_{n-1}^{'}}{p_{n-1}^{'}}.$$
即有
    \begin{align*}
    p_n&=a_0p'_{n-1}+q'_{n-1},\\
    q_n&=p'_{n-1}.
    \end{align*}
\par
将归纳假设用于$[a_1;a_2,\cdots,a_n]$的收敛因子,有
    \begin{align*}
    p'_{n-1}=a_np'_{n-2}+p'_{n-3},\\
    q'_{n-1}=a_nq'_{n-2}+q'_{n-3},
    \end{align*}
从而有
    \begin{align*}
    p_n&=a_0(a_np'_{n-2}+p'_{n-3})+(a_nq'_{n-2}+q'_{n-3}) \\
    &=a_n(a_0p'_{n-2}+q'_{n-2})+(a_0p'_{n-3}+q'_{n-3}) \\
    &=a_np_{n-1}+p_{n-2},\\
    q_n&=a_np'_{n-2}+p'_{n-3}=a_nq_{n-1}+q_{n-2}.
    \end{align*}
由归纳法知命题得证.\qed
\par
\textbf{推论 2.1.2.  }\textsuperscript{\cite{Khinchin}}
$\forall k\geqslant 1$,都有
$$q_kp_{k-1}-p_kq_{k-1}=(-1)^k.$$
或写为
$$\frac{p_{k-1}}{q_{k-1}}-\frac{p_k}{q_k}=\frac{(-1)^k}{q_kq_{k-1}}.$$
\par
\textbf{推论 2.1.3.  }\textsuperscript{\cite{Khinchin}}
$\forall k\geqslant2$,都有
$$q_kp_{k-2}-p_kq_{k-2}={(-1)}^{k-1}a_k.$$
或写为
$$\frac{p_{k-2}}{q_{k-2}}-\frac{p_k}{q_k}=\frac{(-1)^{k-1}a_k}{q_kq_{k-2}}.$$
\par
根据\textbf{推论  2.1.2},\textbf{推论  2.1.3},我们不难得到:\par
\textbf{推论 2.1.4.  }\textsuperscript{\cite{Khinchin}}
偶数阶收敛因子单调递增,奇数阶收敛因子单调递减,且所有偶数项收敛因子都小于奇数项收敛因子.
\par
\textbf{命题 2.1.5.  }\textsuperscript{\cite{Khinchin}}\textbf{(连分数收敛性的判定定理).  }
无限连分数$[a_0;a_1,a_2,\cdots,a_n,\cdots]$收敛的充分必要条件是级数
\begin{equation}\label{series}
\sum_{n=1}^{\infty}a_n
\end{equation}
发散.
\par
\textbf{证明:  }
由\textbf{推论 2.1.4}可以知道, $[a_0;a_1,a_2,\cdots,a_n,\cdots]$的奇、偶数项收敛因子都是收敛的.故$[a_0;a_1,a_2,\cdots,a_n,\cdots]$收敛等价于其奇、偶数项收敛因子收敛于同一极限.根据(\ref{rfc}),知道这又等价于
$$\lim_{k\to\infty}{q_kq_{k+1}}=+\infty.$$\par
\textbf{必要性: }假设(\ref{series})收敛,由(\ref{rfc})可知
$$q_k>q_{k-2},\quad \forall k\geqslant 2.$$
因此, $q_k>q_{k-1}$和$q_{k-1}>q_{k-2}$至少有一个成立.\par
若前者成立,则由(\ref{rfc}),有
$$q_k=a_kq_{k-1}+q_{k-2}<a_kq_k+q_{k-2}.$$
由级数的收敛性知,存在正整数$k_0$,使得$\forall k\geqslant k_0$,都有$a_k<1$,从而
$$q_k<\frac{q_{k-2}}{1-a_k}.$$\par
若后者成立,则
$$q_k=a_kq_{k-1}+q_{k-2}<(1+a_k)q_{k-1}<\frac{q_{k-1}}{1-a_k}.$$\par
故无论如何,对足够大的$k$总存在$l<k$,使得
$q_k<\dfrac{1}{1-a_k}q_l.$
我们重复这样的步骤,得到
$$q_k<\frac{q_s}{\left(1-a_k\right)\left(1-a_{k_1}\right)\ldots\left(1-a_{k_m}\right)},$$
其中$k>k_1>\cdots>k_m\geqslant k_0>s$.然而,由于(\ref{series})收敛,故无穷乘积
$$\prod_{n=k_0}^{\infty}\left(1-a_n\right)$$
有确定的正值,记为$\xi$,并且显然有
\begin{equation}\label{repeated}
(1-a_k)(1-a_{k_1})\cdots(1-a_{k_m})\geqslant \prod_{n=k_0}^{\infty}(1-a_n)=\lambda.   
\end{equation}
我们再记$q_0,q_1,\cdots,q_{k_0-1}$中最大的为$M$.
由(\ref{repeated})可知, $\forall k\geqslant k_0$,有
$$q_k<\frac{M}{\xi},$$
从而
$$q_{k+1}q_k<\frac{M^2}{\lambda^2}<\infty.$$
这就证明了连分数的收敛.\par
\textbf{充分性: }假设(\ref{series})发散.记$q_0,q_1$中较小的为$m>0$.由$q_k>q_{k-2},\forall k\geqslant2$,可知总有$q_k>m$.因而
$$q_k\geqslant q_{k-2}+ma_k,\quad \forall k\geqslant 2,$$
从而
\begin{align*}
q_{2k} &\geqslant q_0+c\sum_{n=1}^k a_{2n},\\
q_{2k+1} &\geqslant q_1+c\sum_{n=1}^k a_{2n+1}.
\end{align*}
上面两式相加可得
$$q_{2k}+q_{2k+1}>q_{0}+q_{1}+c\sum_{i=1}^{2k+1}a_{n},$$
也就是说
$$q_k+q_{k-1}>c\sum\limits_{n=1}^k a_n.$$
又由均值不等式可知
$$q_{k}q_{k-1}>\frac{c^{2}}{2}\sum_{n=1}^{k}a_{n}.$$
结合级数(\ref{series})的发散,说明了连分数的收敛.\qed
\par
\textbf{推论 2.1.6.  }\textsuperscript{\cite{Khinchin}}
正规无限连分数总是收敛的.
\par
在本文的以下内容中,我们只考虑正规连分数,并简称为连分数.称无限连分数的极限值为连分数的值.\par
不难验证,有限连分数$[a_0;a_1,a_2,\cdots,a_k,1]=[a_0;a_1,a_2,\cdots,a_k+1]$,因此我们可以将最后一个部分商为$1$的有限连分数排除在我们考虑的范围之外.在现在的讨论范围内,我们有以下的重要定理.\par
\textbf{命题 2.1.7.  }\textsuperscript{\cite{Khinchin}}
$\forall x\in \mathbb{R}$,都有唯一的一个连分数的值等于$x$,并且该连分数是无限的当且仅当$x$是无理的.
\par
我们将不对这个定理进行证明,实际上,在介绍了Gauss映射$\tau$之后,实数与连分数的对应关系就被建立起来了.\par
在最后,我们对收敛因子$\dfrac{p_k}{q_k}$中分母$q_k$的大小进行估计.\par
\textbf{命题 2.1.8.  }\textsuperscript{\cite{Khinchin}}
$\forall k\geqslant 1$,都有
$$q_{k}\geqslant2^{\frac{k-1}{2}}.$$
\par
\textbf{证明:  }
由(\ref{rfc}),可知
$$q_k=a_k q_{k-1}+q_{k-2}\geqslant q_{k-1}+q_{k-2}\geqslant2q_{k-2}.$$
并且$q_0=1,q_1=a_1\geqslant 1$,可知命题成立.\qed
\par
结合(\ref{rfc}),可以得到下面估计
\begin{equation}\label{klength}
\left|\frac{p_n}{q_n}-\frac{p_n+p_{n-1}}{q_n+q_{n-1}}\right|=\frac{1}{q_n(q_n+q_{n-1})}\leqslant\frac{1}{2^{n-1}}.
\end{equation}
这个式子在以后会用到.
\subsection{连分数度量理论简介}
在介绍连分数的度量理论之前,我们先不加证明地指出下面的四个结论.\par
\textbf{命题 2.2.1.  }\textsuperscript{\cite{Khinchin}}
对任意的正函数$\varphi(q)$,总存在无理数$\alpha$使得不等式
$$\left|\alpha-\frac{p}{q}\right|<\varphi\left(q\right)$$
对无穷多对整数$(p,q),q>0$成立.
\par
\textbf{命题 2.2.2.  }\textsuperscript{\cite{Khinchin}}
记无理数$\alpha$的连分数表示为$[a_0;a_1,a_2,\cdots,a_n,\cdots]$.若$\{a_n\}_{n\in \mathbb{N}}$有界,则存在足够小的$c>0$,使得不等式
$$\left|a-\frac{p}{q}\right|<\frac{c}{q^2}$$
对任意整数对$(p,q),q>0$都不成立.\par
反过来,若$\{a_n\}_{n\in \mathbb{N}}$无界,则存在无穷多的整数对$(p,q),q>0$使得上面的不等式成立.
\par
\textbf{命题 2.2.3.  }\textsuperscript{\cite{Khinchin}}\textbf{(Liouville定理).   }
对$n$阶无理代数数$\alpha$,存在正数$C$使得对任意整数对$(p,q),q>0$,都成立不等式
$$\left|\alpha-\frac{p}{q}\right|>\frac{C}{q^n}.$$
\par
(其中, $\alpha$是$n$阶代数数指的是它满足一个$n$阶整系数多项式方程,而不满足任何$k<n$阶整系数多项式方程.不是代数数的数称为超越数.)\par
关于超越数,作为Liouville定理的推论,有\par
\textbf{推论 2.2.4.  }\textsuperscript{\cite{Khinchin}}
若对任意$C>0$与任意的正整数$n$,都存在整数对$(p,q),q>0$使得
$$\left|\alpha-\frac pq\right|\leqslant\frac C{q^n},$$
则$\alpha$是超越数.
\par
根据\textbf{命题 2.2.2}与\textbf{命题 2.2.3},我们还可以知道:任意的二阶实代数数,其连分数表示中的全体部分商有界.\par
从以上命题不难发现,实数的许多算术性质同其连分数表示形式有比较紧密的联系.于是可以提出一个自然的问题:对于一个具体性质而言,满足性质的和不满足性质的数,那些更“普遍”?为解决这样的问题,需要建立连分数的度量理论.\par
$\forall x \in\mathbb{R}$, $x$存在唯一的连分数表示
$$x=[a_0;a_1,a_2,\dots],$$
从而,也建立起了部分商作为实数$x$的函数$a_n(x)=a_n$.
对于有理数$x$,它被表示为有限连分数,在这里,对于尚没有定义的$a_n(x)$,可以定义$a_n(x)=\infty$.但如果只考虑关于度量的性质的话,可以忽略$x$为有理数的情况,因为在Lebesgue测度$\lambda$和Gauss测度$\gamma$考虑的度量意义下,有理数是零测的.\par
为使所考虑的测度是有限的,可以通过限定$x\in[0,1)$,此时$a_0(x)=0.$\par
先考虑$n=1,2$的简单情况,这位我们认识一般的$a_n(x)$有很大作用.\par
由定义,
$$a_1(x)=k\iff k\leqslant x<k+1\iff \frac{1}{k+1}<x\leqslant\frac{1}{k}.$$
故可以根据$a_1(x)$的值将$[0,1)$划分为一些区间
$$[\frac{1}{k+1},\frac{1}{k}),\quad k\in \mathbb{N^{+}}.$$
这些区间都是某个一阶柱集.\par
同样地,
$$a_2=l\iff\exists k,a_1=k,a_2=l\iff\exists k,\dfrac{1}{k+\frac{1}{l}}\leqslant x<\dfrac{1}{k}+\frac{1}{l+1}.$$
故可以根据$a_2(x)$将一阶柱集$[\frac{1}{k+1},\frac{1}{k})$划分为一些区间
$$[\dfrac{1}{k+\frac{1}{l}},\dfrac{1}{k+\frac{1}{l+1}}),\quad l\in \mathbb{N^{+}}.$$
这些区间都是某个二阶柱集.\par
总的来说, $a_{k}(x)$将一个$k-1$阶柱集划分为可数个$k$阶柱集,在每个$k$阶柱集上, $a_{k}(x)$有恒定的值.并且当$k$为奇数时, $a_k(x)$在每个柱集上的值,从右到左递增;当$k$为偶数时, $a_k(x)$在每个柱集上的值,从左到右递增.并且有下面的定量关系:\par
\textbf{命题 2.2.5.  }\textsuperscript{\cite{Khinchin}}
给定一组正整数$(i_1,i_2,\cdots,i_k)$,它定义了一个$k$阶柱集,其端点为由$(a_1=i_1,a_2=i_2,\cdots,a_k=i_k)$所决定的
$$\frac{p_k+p_{k-1}}{q_k+q_{k-1}}\quad\text{和}\quad\frac{p_k}{q_k}.$$
\par
对一般的$E=\{x\in [0,1)|a_{n_1}=i_1,a_{n_2}=i_2,\cdots,a_{n_k}=i_k\}$,有递推关系:
\begin{align*}
\sum_{i_{\ell}=1}^{\infty}&\{x\in [0,1)|a_{n_1}=i_1,\cdots,a_{n_{\ell-1}}=i_{\ell-1},a_{n_{\ell}}=i_{\ell},a_{n_{\ell+1}}=i_{\ell+1},\cdots,a_{n_k}=i_k\}\\
&=\{x\in [0,1)|a_{n_1}=i_1,\cdots,a_{n_{\ell-1}}=i_{\ell-1},a_{n_{\ell+1}}=i_{\ell+1},\cdots,a_{n_k}=i_k\}
\end{align*}
成立,其中对集合的$\sum$表示集合的无交并.\par
为呼应本节开头的内容,并举例说明连分数的度量性质在研究连分数的算数性质上的重要性,我们不加证明地介绍下面的两个性质并结束这一节.\par
\textbf{命题 2.2.6.  }\textsuperscript{\cite{Khinchin}}
在连分数表示下,部分商$\{a_n\}_{n\in \mathbb{N}}$有界的点集是一个Lebesgue零测集.
\par
\par
\textbf{命题 2.2.7.  }\textsuperscript{\cite{Khinchin}}
设$f$是定义在正半轴上的连续函数,且$xf(x)$不递增,则当积分
$$
\int\limits_{c}^{\infty}f\left(x\right)dx
$$
对某个正数$c$发散时,不等式
$$\left|\alpha-\frac{p}{q}\right|<\frac{f(q)}{q}$$
对几乎处处的$\alpha$,有无穷多的整数对对$(p,q),q>0$成立.
\par

\subsection{连分数与Gauss变换}
在这一节中,我们将定义Gauss变换$\tau$和Gauss测度$\gamma$.\par
\textbf{定义 2.3.1.  }
记$I=[0,1)$, $\mathcal{B}_I$为$I$的Borel $\sigma$代数,定义Gauss变换$\tau:I\to I:$
$$
\tau\left(x\right)=\left\{
\begin{array}{cc}
x^{-1}-\left\lfloor x^{-1}\right\rfloor&\text{当}x\neq0,\\
0&\text{当}x=0.\end{array}\right.
$$
\par
可以知道, $I$中实数和其连分数部分商的关系为
$$a_1=a_1(x)=\lfloor1/x\rfloor,\quad a_n=a_n(x)=\lfloor1/\tau^{n-1}(x)\rfloor,$$
以及
\begin{equation}\label{tauandinquo}
a_n\left(x\right)=a_1\left(\tau^{n-1}\left(x\right)\right).
\end{equation}
(其中,定义$a_1(0)=\infty$.)\par
历史上, Gauss问题是连分数度量理论中第一个问题,其内容是:估计误差
$$m_n(x)-\frac{\ln(1+x)}{\ln{2}}$$
其中, $m_n(x)=\lambda({u\in[0,1)\mid\tau^n(u)\leqslant x})$, $\lambda$为Lebesgue测度.这个问题促使了Gauss测度的诞生.\par
\textbf{定义 2.3.2.  }
Gauss测度$\gamma:\mathcal{B}_I\to[0,1]:$
$$\gamma\left(A\right)=\frac{1}{\ln2}\int_A\frac{\mathrm{d}x}{x+1},\quad A\in\mathcal{B}_I.$$
\par
下面给出Gauss测度与Lebesgue测度的两个性质.\par
\textbf{命题 2.3.1.  }\textsuperscript{\cite{Iosifescu}}
Gauss变换$\tau$是保Gauss测度$\gamma$的而不保持Lebesgue测度$\lambda$.
\par
\textbf{证明:  }
前半部分,即证
\begin{equation}\label{Gausspreserved}
\gamma\left(\tau^{-1}(A)\right)=\gamma(A),\quad \forall A\in\mathcal{B}_I.
\end{equation}
这里只对$A=[0,u)$验证(\ref{Gausspreserved}).至于从$A=[0,u)$的情况过渡到$\forall A\in\mathcal{B}_I$的情况,之后会进行介绍.\par
考虑$A$
在Gauss变换之逆$\tau^{-1}$下的原像,有
$$\tau^{-1}\left([0,u)\right)=\bigcup_{i\in\mathbf{N}_+}\left(\frac{1}{u+i},\frac{1}{i}\right],$$
故(\ref{Gausspreserved})即
$$\int_0^u\frac{\mathrm{d}x}{x+1}=\sum\limits_{i\in\mathbf{N}_+}\int_{1/(u+i)}^{1/i}\frac{\mathrm{d}x}{x+1}.$$
上式直接积分可验证是成立的.
\par
后半部分,为证Gauss变换$\tau$不保持Lebesgue测度$\lambda$,只需找到一个反例,取$A=(\dfrac{1}{2},1)$,有
$$
\begin{aligned}
\lambda\left(\tau^{-1}\left(A\right)\right)& =\sum_{i\in\mathbf{N}_{+}}\left(\frac{1}{i+1/2}-\frac{1}{i+1}\right)=2\sum_{i\in\mathbf{N}_{+}}\left(\frac{1}{2i+1}-\frac{1}{2i+2}\right)  \\
&=2\left(\log2-1+\frac12\right)=2\log2-1.
\end{aligned}
$$
而
$$\lambda(A)=\frac{1}{2}\neq2\log2-1=\lambda\left(\tau^{-1}\left(A\right)\right).$$
这就完成了证明.\qed
\par
\textbf{命题 2.3.2.  }\textsuperscript{\cite{Liu_Peng}}
Gauss测度$\gamma$和Lebesgue测度$\lambda$等价.
\par
\textbf{证明:  }
由于
$$\frac{1}{2\ln2}\leqslant \frac{1}{(\ln2)(x+1)}\leqslant \frac{1}{\ln2},\quad \forall x\in I.$$
故$\forall A\in \mathcal{B_I}$
$$\frac{\lambda(A)}{2\ln2}=\int_A \frac{1}{2\ln2}\mathrm{d}x\leqslant \int_A\frac{1}{(\ln2)(x+1)}\mathrm{d}x=\gamma(A)\leqslant \int_A\frac{1}{\ln2}\mathrm{d}x=\frac{\lambda(A)}{\ln2}.$$
所以两个测度等价.\qed
\par
作为本节的末尾,我们介绍由Gauss变换$\tau$启发的另一个变换$\bar{\tau}:\Omega^2\to\Omega^2,$
$$\bar{\tau}\left(\omega,\theta\right)=\left(\tau\left(\omega\right),\frac{1}{a_1\left(\omega\right)+\theta}\right),$$
其中$\Omega=[0,1]\setminus\mathbb{Q}$.
它是一个双射(注意到$\tau$不是,如$\tau(\dfrac{2}{3})=\tau(\dfrac{2}{5})=\dfrac{1}{2}$),其逆为
$$\bar{\tau}^{-1}\left(\omega,\theta\right)=\left(\frac{1}{a_1\left(\theta\right)+\omega},\tau\left(\theta\right)\right),$$
并且有
$$\bar{\tau}^{n}\left(\omega,\theta\right)=\left(\tau^{n}\left(\omega\right),\left[a_{n}\left(\omega\right),\cdots,a_{2}\left(\omega\right),a_{1}\left(\omega\right)+\theta\right]\right),$$
和
$$\bar{\tau}^{-n}\left(\omega,\theta\right)=\left([a_{n}\left(\theta\right),\cdots,a_{2}\left(\theta\right),a_{1}\left(\theta\right)+\omega],\tau^{n}\left(\theta\right)\right).$$\par
除了可逆外, $\bar\tau$的另一个优点是, (与$\tau$的情况下定义$\{a_{n}\}_{n\in\mathbb{N}}$类似)我们可以定义$\{\bar{a}_{\ell}\}_{\ell\in\mathbb{Z}}$,为
$$\bar{a}_{\ell+1}\left(\omega,\theta\right)=\bar{a}_1\left(\bar{\tau}^\ell\left(\omega,\theta\right)\right).$$
两者之间的关系是
\begin{equation}\label{taubartauan}
\bar{a}_{0}\left(\omega,\theta\right)=a_{1}\left(\theta\right),\bar{a}_{n}\left(\omega,\theta\right)=a_{n}\left(\omega\right),\bar{a}_{-n}\left(\omega,\theta\right)=a_{n+1}\left(\theta\right),\forall n\in\mathbb{N}.
\end{equation}
这一点在之后会用到.\par
$\bar\tau$也引出一种Gauss测度$\bar\gamma:\mathcal{B}_{\Omega^2}\to[0,1]$
$$\bar{\gamma}(A)=\frac{1}{\log2}\iint_A\frac{\mathrm{d}x\mathrm{d}y}{\left(xy+1\right)^2}.$$
与$\gamma$的关系是
$$\bar{\gamma}(A\times I) = \bar{\gamma}(I\times A) = \gamma(A)$$
并且也有与\textbf{命题 2.3.2.}类似的结论.\par
\textbf{命题 2.3.3.  }\textsuperscript{\cite{Iosifescu}}
$\bar\tau$是保$\bar\gamma$测度的.
\par
\subsection{部分商序列的性质}
在Gauss测度$\gamma$的意义下,连分数部分商序列有以下性质:\par
\textbf{命题 2.4.1.  }\textsuperscript{\cite{Liu_Peng,Iosifescu}}
部分商序列是严平稳的,从而是同分布的.
\par
\textbf{证明:  }
由(\ref{tauandinquo})与\textbf{命题 2.3.1}即可证明.\qed
\par
\textbf{命题 2.4.2.  }\textsuperscript{\cite{Liu_Peng,Iosifescu}}
连分数序列可逆,即
$$a_\ell:m\leqslant\ell\leqslant n\quad \text{和}\quad a_{m+n-\ell}:m\leqslant\ell\leqslant n$$
同分布.
\par
\textbf{证明:  }
类似\textbf{命题 2.4.1},可以证明对$\bar\tau$,其定义的$\{\bar{a}_n\}_{n\in\mathbb{Z}}$也是严平稳的(在$\bar\gamma$测度下).故
$${\bar{a}}_\ell:m\leqslant \ell\leqslant n\quad \text{和}\quad {\bar{a}}_{\ell-m-n+1}:m\leqslant \ell\leqslant n$$
同分布.\par
结合(\ref{taubartauan})知,
$$a_\ell:m\leqslant\ell\leqslant n\quad \text{和}\quad a_{m+n-\ell}:m\leqslant\ell\leqslant n$$
同分布.\qed
\par
\textbf{命题 2.4.3.  }\textsuperscript{\cite{Liu_Peng}}
连分数部分商序列不独立.
\par
\textbf{证明:  }
只需要举出反例,下面说明$\gamma(a_1=1,a_2=1)\neq\gamma(a_1=1)\gamma(a_2=1)$.\par
由\textbf{命题 2.3.1.}可知
$$\gamma(a_1=1,a_2=1)=\frac{1}{\ln2}\int^{\frac{2}{3}}_\frac{1}{2}\frac{1}{1+x}\mathrm{d}x=\frac{1}{\ln2}\ln\frac{10}{9};$$
$$\gamma(a_1=1)=\frac{1}{\ln2}\int^1_\frac{1}{2}\frac{1}{1+x}\mathrm{d}x=\frac{1}{\ln2}\ln\frac{4}{3};$$
$$\gamma(a_2=1)=\gamma(a_1=1)=\frac{1}{\ln2}\ln\frac{4}{3}.$$
最后一个式子是因为部分商序列同分布.\par
从而$\gamma(a_1=1,a_2=1)\neq\gamma(a_1=1)\gamma(a_2=1)$,连分数部分商序列不独立.\qed


\sectionbreak
% ---------------------------------------------------------------------------- %
\section{连分数动力系统的$\psi$-混合性}

\subsection{保测系统简介}
\textbf{定义 3.1.1.  }
记两个概率空间为$(X_i,\mathcal{B}_i,\mu_i),i=1,2$,映射$T:X_1\to X_2$被称为可测的,若
$$\forall A_2\in\mathcal{B}_2,\quad \exists A_1\in\mathcal{B}_1,\quad T^{-1}(A_2)=A_1;$$
被称为保测的,若
$$\mu_1(T^{-1}(A_2))=\mu_2(A_2).$$
若变换$T:X\to X$是保测的,则称$(X,\mathcal{B},\mu,T)$为保测系统.
\par
下面我们将介绍保测变换的验证方法,以解决\textbf{命题 2.3.1}的证明中所保留的,从$A=[0,u)$的情况过渡到$A\in\mathcal{B}_I$的情况.在此之前,一些准备工作是需要的.\par
\textbf{定义 3.1.2.  }
记$X$为一个集合, $\mathcal{S}$是$X$的一个子集类.若有:
\begin{enumerate}
    \item $\emptyset\in S$;
    \item 若$A,B\in S$,则$A\cap B\in S$;
    \item 若$A\in S$,则$\exists E_i\in \mathcal{S},i=1,2,\cdots,n$使得$X\setminus A=\sum\limits_{i=1}^n E_i$.
\end{enumerate}
则称$S$为一个半代数.\par
\par
\textbf{定义 3.1.3.  }
记$X$为一个集合, $\mathcal{S}$是$X$的一个子集类.若有:
\begin{enumerate}
    \item $\emptyset\in \mathcal{A}$;
    \item 若$A,B\in \mathcal{A}$,则$A\cap B\in \mathcal{A}$;
    \item 若$A\in \mathcal{A}$,则$X\setminus \mathcal{A}\in S$.
\end{enumerate}
则称$\mathcal{A}$为一个代数.\par
\par
不难验证,代数是半代数, $\sigma$代数是代数.代数的交仍是代数,并且对$X$的任意子集族$\mathcal{G}$, $2^X$是一个包含$\mathcal{G}$的半代数,从而也是代数, $\sigma$代数.\par
由上面的讨论,可知以下三个子集族的定义是合理的.\par
\textbf{定义 3.1.4.  }
\begin{enumerate}
    \item $\mathcal{A}(\mathcal{S})$为包含半代数$\mathcal{S}$的最小代数,也称由半代数$\mathcal{S}$生成的代数,
    \item $\mathcal{B}(\mathcal{A})$为包含代数$\mathcal{A}$的最小$\sigma$代数,也称由代数$\mathcal{A}$生成的$\sigma$代数,
    \item $\mathcal{B}(\mathcal{A}(\mathcal{S}))$称由半代数$\mathcal{S}$生成的$\sigma$代数.
\end{enumerate}
\par
\textbf{定义 3.1.5.  }
记$X$为一个集合,$\mathcal{C}$是$X$的一个子集类.若有:
\begin{enumerate}
    \item $\mathcal{C}$对单调增集合列之并封闭,即
          $$\text{若}E_{1}\subset E_{2}\subset\cdots,E_{i}\in{\mathcal C}\text{,则}\bigcup_{i=1}^{\infty}E_{i}\in{\mathcal C};$$
    \item $\mathcal{C}$对单调减集合列之交封闭,即
          $$\text{若}F_{1}\supset F_{2}\supset\cdots,F_{i}\in{\mathcal C}\text{,则}\bigcap_{i=1}^{\infty}F_{i}\in{\mathcal C}.$$
\end{enumerate}
则称$\mathcal{C}$为单调类.
\par
不难验证,单调类的交仍是单调类,并且对$X$的任意子集族$\mathcal{G}$, $2^X$是包含$\mathcal{G}$的单调类.\par
由上面的讨论,可知以下子集族的定义是合理的.\par
\textbf{定义 3.1.6.  }$\mathcal{C}(\mathcal{G})$为包含$\mathcal{G}$的最小单调类,也称由$\mathcal{G}$生成的单调类.
\par
下面将介绍保测变换的一个验证方法,先介绍两个引理:\par
\textbf{引理 3.1.1.  }\textsuperscript{\cite{Ergodic_Sun,Ergodic_theory}}
由半代数$\mathcal{S}$生成的代数为:
\begin{equation}\label{genealge}
    {\mathcal A}({\mathcal S})=\left\{E=\sum\limits_{i=1}^{n}E_{i}|E_{i}\in{\mathcal S},i=1,2,\cdots,n,n\in\mathbb{N^{+}}\right\}.
\end{equation}
\par
\textbf{证明:  }
先记(\ref{genealge})的右边为$\mathcal{G}$.\par
首先,根据定义,验证$\mathcal{G}$为一个代数.代数定义的第1条显然成立.\par
对$\mathcal{G}$中的两个集合
$$E=\sum\limits_{i=1}^{n}E_{i},E_{i}\in{\mathcal S},i=1,2,\cdots,n,\quad F=\sum\limits_{j=1}^{m}F_{j},F_{j}\in{\mathcal S},j=1,2,\cdots,m,$$
可知
$$E\cap F=\bigcup\limits_{(i,j)\in\{1,2,\cdots,n\}\times\{1,2,\cdots,m\}}E_i\cap F_j.$$
由于$\mathcal{S}$是半代数,故诸$E_i\cap F_j\in\mathcal{S}$.从而$E\cap F\in \mathcal{G}.$即代数定义的第2条成立.\par
而对
$$X\setminus E=X\setminus(\sum\limits_{i=1}^n E_i)=\bigcap\limits_{i=1}^n (X\setminus E_i).$$
由于$E_i \in \mathcal{S}$,且$\mathcal{S}$为半代数,故
$$X\setminus{E_i}=\sum\limits_{j=1}^{m_i}A_{ij},\quad i=1,2,\cdots,n,j=1,2,\cdots,m_i,$$
从而$X\setminus E_i\in \mathcal{G}$.又由已经证明的定义中第2条,可知$X\setminus E=\bigcap\limits_{i=1}^n (X\setminus E_i)\in \mathcal{G}$.从而定义中第3条也成立.综上所述, $\mathcal{G}$是一个代数.\par
进一步地,再证明$\mathcal{G}$是包含$\mathcal{S}$的最小代数.任取包含$\mathcal{S}$的代数$\mathcal{A}$和$E\in \mathcal{G},E=\sum\limits_{i=1}^{n}E_{i},E_{i}\in{\mathcal S}$.由$E_i\in \mathcal{S}\subset \mathcal{A}$,有$X\setminus E_i\in \mathcal{A}$.由$X\setminus E=\bigcap\limits_{i=1}^n (X\setminus E_i)$知$X\setminus E\in \mathcal{A}$.从而
$$X\setminus (X\setminus E)=E\in\mathcal{A}.$$
这就证明了$\mathcal{G}\subset\mathcal{A}$,由$\mathcal{A}$的任意性,知$\mathcal{G}$是包含$\mathcal{S}$的最小代数.\qed
\par
\textbf{引理 3.1.2.  }\textsuperscript{\cite{Ergodic_Sun,Ergodic_theory}}
由代数$\mathcal{A}$生成的$\sigma$代数为:
$$\mathcal{B}(\mathcal{A})=\mathcal{C}(\mathcal{A}).$$
\par
\textbf{证明:  }
我们先证明$\mathcal{C}(\mathcal{A})$对取补运算封闭.记$\mathcal{C}_c(\mathcal{A})$是$\mathcal{C}(\mathcal{A})$的子集,且满足:
$$\text{若}A\in\mathcal{C}_c(\mathcal{A}),\text{则}A^{c}\in\mathcal{C}_c(\mathcal{A}).$$\par
由于$\mathcal{A}\subset\mathcal{C}_c(\mathcal{A})$知$\mathcal{C}_c(\mathcal{A})$非空.设$\{E_i\}_{i\in\mathbb{N}^+}$是$\mathcal{C}_c(\mathcal{A})$中的单调增集合列.则$\{(E_i)^{c}\}_{i\in\mathbb{N}}$是$\mathcal{C}_c(\mathcal{A})$中的单调减集合列,自然它们都是$\mathcal{C}(\mathcal{A})$中的单调集合列,故
$$E=\bigcup_{i=1}^{\infty}E_{i}\in{\mathcal C}\left({\mathcal A}\right),E^c=\bigcap_{i=1}^{\infty}E_{i}^{c}\in{\mathcal C}\left({A}\right).$$
从而$E\in\mathcal{C}_c(\mathcal{A})$, $\mathcal{C}_c(\mathcal{A})$对单调增集合列之并封闭.\par
运用相同的办法,可以证明$\mathcal{C}_c(\mathcal{A})$对单调减集合列之交封闭.从而$\mathcal{C}_c(\mathcal{A})$是单调类, $\mathcal{C}(\mathcal{A})\subset\mathcal{C}_c(\mathcal{A})$.故$\mathcal{C}_c(\mathcal{A})=\mathcal{C}(\mathcal{A})$,从而$\mathcal{C}(\mathcal{A})$对取补运算封闭.\par
再证明$\mathcal{C}(\mathcal{A})$对有限交封闭.任取$F\in\mathcal{C}(\mathcal{A})$.记$\mathcal{C}_F(\mathcal{A})$是$\mathcal{C}(\mathcal{A})$的子集,且满足
$$\text{若}E\cap F\in\mathcal{C}(\mathcal{A}),\text{则}E\in\mathcal{C}_F(\mathcal{A}).$$\par
由于$\emptyset\in\mathcal{C}_F(\mathcal{A})$,知$\mathcal{C}_F(\mathcal{A})$非空.设$\{E_i\}_{i\in\mathbb{N}^+}$是$\mathcal{C}(\mathcal{A})$中的单调增集合列.则$\{E_i\cap F\}_{i\in\mathbb{N}^+}$是$\mathcal{C}_F(\mathcal{A})$中的单调增集合列,自然也是$\mathcal{C}(\mathcal{A})$中的单调增集合列,故
$$E\cap F=\bigcup_{i=1}^{\infty}\left(E_i\cap F\right)\in\mathcal{C}\left(A\right)$$
从而$E\in\mathcal{C}_F(\mathcal{A})$. $\mathcal{C}_F(\mathcal{A})$对单调增集合列之并封闭.\par
运用相同的办法,可以证明$\mathcal{C}_F(\mathcal{A})$对单调减集合列之交封闭.从而$\mathcal{C}_F(\mathcal{A})$是单调类, $\mathcal{C}(\mathcal{A})\subset\mathcal{C}_F(\mathcal{A})$.故$\mathcal{C}_F(\mathcal{A})=\mathcal{C}(\mathcal{A})$, $\mathcal{C}$对交$F$封闭.又由$F$的任意性可以证得,$\mathcal{C}(\mathcal{A})$对有限交运算封闭.\par
故$\mathcal{C}(\mathcal{A})$是代数,对有限的交(并)封闭.又由$\mathcal{C}(\mathcal{A})$是单调类,可知其对无限的交(并)封闭,从而是$\sigma$代数.\qed
\par
\textbf{命题 3.1.3.  }\textsuperscript{\cite{Ergodic_Sun,Ergodic_theory}}\textbf{(保测变换的验证方法).  }
$(X,\mathcal{B},\mu)$为一概率空间, $T:X\to X$为可测变换. $\mathcal{B}$为半代数$\mathcal{S}$生成的$\sigma$代数.若
$$\mu(T^{-1}(A))=\mu(A),\quad \forall A\in \mathcal{S},$$
成立,则$T$为保测的.
\par
\textbf{证明:  }
记$\{A_i\}_{i\in\mathbb{N^+}}$中的任意不交集合列,由\textbf{引理 3.1.1}, $\mathcal{A}(\mathcal{S})$中的元素都有$\sum\limits_{i=1}^n A_i$的形式.由于在$\mathcal{S}$上满足$\mu(T^{-1}(A))=\mu(A)$,结合
$$\mu(T^{-1}(\sum\limits_{i=1}^n A_i))=\mu(\sum\limits_{i=1}^n T^{-1}(A_i))=\sum\limits_{i=1}^n\mu(T^{-1}(A_i))=\sum\limits_{i=1}^n\mu(A_i)=\mu(\sum\limits_{i=1}^n A_i),$$
知$\mu(T^{-1}(A))=\mu(A)$在$\mathcal{A}(\mathcal{S})$上也成立.\par
记$\{B_n\}_{n\in\mathbb{N}^+}$是满足$\mu(T^{-1}(B_i))=\mu(B_i)$的单调增集合列,由于
$$\mu\left(T^{-1}\bigcup_{i=1}^{n}B_{i}\right)=\mu\left(T^{-1}B_{n}\right)=
    \mu\left(B_{n}\right)=\mu\left(\bigcup_{i=1}^{n}B_{i}\right).$$
并且$\mu\left(\bigcup\limits_{i=1}^{n}B_{i}\right),\mu\left(T^{-1}\bigcup\limits_{i=1}^{n}B_{i}\right)$为单调有界数列,是收敛的,从而
\begin{align*}
    \lim_{n\to\infty}\mu\left(\bigcup\limits_{i=1}^{n}B_{i}\right)       & = \mu\left(\bigcup\limits_{i=1}^{\infty}B_{i}\right),       \\
    \lim_{n\to\infty}\mu\left(T^{-1}\bigcup\limits_{i=1}^{n}B_{i}\right) & = \mu\left(T^{-1}\bigcup\limits_{i=1}^{\infty}B_{i}\right). \\
    \mu\left(T^{-1}\bigcup_{i=1}^{\infty}B_{i}\right)                    & =\mu\left(\bigcup_{i=1}^{\infty}B_{i}\right).
\end{align*}
可知满足$\mu(T^{-1}(A))=\mu(A)$的子集类是包含$\mathcal{A}(\mathcal{S})$的单调类.结合\textbf{引理 3.1.2}知,满足$\mu(T^{-1}(A))=\mu(A)$的子集类是$\mathcal{B}(\mathcal{A}(\mathcal{S}))=\mathcal{B}$.\qed
\par
凭借\textbf{命题 3.1.3}可以为\textbf{命题 2.3.1}给出完整的证明.
\subsection{混合性简介}
\textbf{定义 3.2.1.  }
若保测系统$(X,\mathcal{B},\mu,T)$满足
$$\lim_{n\to\infty}{\frac{1}{n}}\sum_{i=0}^{n-1}|\mu(T^{-i}A\cap B)-\mu(A){\mu}(B)|=0$$
对任意可测集$A,B\in\mathcal{B}$成立,则称该保测系统为弱混合的.
\par
\textbf{定义 3.2.2.  }
若保测系统$(X,\mathcal{B},\mu,T)$满足
$$\lim\limits_{n\to\infty}\left|\mu(T^{-n}A\cap B)-\mu(A)\mu(B)\right|=0$$
对任意可测集$A,B\in\mathcal{B}$成立,则称该保测系统为强混合的.
\par
记$\{X_k\}_{k\in\mathbb{N^+}}$是定义在概率空间$(X,\mathcal{B},\mu)$上的随机过程,$X_{i}^{j}$为由$X_{\ell},i\leqslant\ell\leqslant j$所生成的$\sigma$代数,其中$i< j,i\in\mathbb{N}^+,j\in\mathbb{N}^+\cup\{\infty\}$.\par
又记$\mathcal{B}_i\subset\mathcal{B},i=1,2$为两个$\sigma$代数,记
$$\alpha(\mathcal{B}_1,\mathcal{B}_2)=\sup_{A\in\mathcal{B}_1,B\in\mathcal{B}_2}\left(|\mu(A\cap B)-\mu(A)\mu(B)|\right);$$
$$\phi(\mathcal{B}_1,\mathcal{B}_2)=\sup_{A\in\mathcal{B}_1,B\in\mathcal{B}_2,\mu(B)>0}\left(|\mu(A|B)-\mu(A)|\right);$$
$$\psi(\mathcal{B}_1,\mathcal{B}_2)=\sup_{A\in\mathcal{B}_1,B\in\mathcal{B}_2,\mu(A)>0,\mu(B)>0}\left(\left|\frac{\mu(A|B)}{\mu(A)}-1\right|\right).$$
\par
\textbf{定义 3.2.3.  }
若随机过程$\{X_k\}_{k\in\mathbb{N^+}}$是严平稳的,并且
$$\lim_{n\to\infty}\alpha(X_{1}^{k},X_{k+n}^{\infty})=0,$$
则称其为$\alpha$-混合的.类似,也可以定义$\phi$-混合、$\psi$-混合.
\par
由于
$$\alpha(\mathcal{A},\mathcal{B})\leqslant\phi(\mathcal{A},\mathcal{B})\leqslant\psi(\mathcal{A},\mathcal{B})$$
故$\psi$-混合性蕴含$\phi$-混合性蕴含$\alpha$-混合性.
\subsection{Kuzmin 定理}
在这一节,我们将证明\textbf{Kuzmin 定理},并由此解决Gauss问题以及为连分数动力系统的混合性证明做铺垫.\par
\textbf{命题 3.3.1.  }\textsuperscript{\cite{Kuzmin,Khinchin,Iosifescu}}\textbf{(Kuzmin定理).  }
记$\{f_n\}_{n\in\mathbb{N}}$是定义在$[0,1]$上的一列实函数,并且满足下面三个条件:
\begin{enumerate}
    \item 递推条件:
          \begin{equation}\label{recur}
              f_{n+1}(x)=\sum\limits_{k=1}^{\infty}\dfrac{1}{(k+x)^{2}}f_{n}\left(\dfrac{1}{k+x}\right);
          \end{equation}
    \item 存在正数$M$,使得$0<f_0(x)<M$;
    \item 存在正数$m$,使得$|f_0'(x)|<m$.
\end{enumerate}
则存在取决于$M,m$的正数$A$与正的常数$\lambda$成立
\begin{equation}\label{Kuzmin}
    f_{n}\left(x\right)=\frac{a}{1+x}+\theta A e^{-\lambda \sqrt{n}},
\end{equation}
$$a=\frac{1}{\ln2}\int_{0}^{1}f_{0}(z)d z,\quad|\theta|<1.$$
\par

该定理的证明较长,在这里我们仅简要叙述主要步骤.先介绍在题设条件下成立的下面四个引理:\par
\textbf{引理 1.  }\textsuperscript{\cite{Iosifescu}}
\begin{equation}\label{lemma1}
    f_{n}\left(x\right)=\sum^{\left(n\right)}f_{0}\left(\frac{p_{n}+x p_{n-1}}{q_{n}+x q_{n-1}}\right)\frac{1}{\left(q_{n}+x q_{n-1}\right)^{2}}
\end{equation}
对任意$n\geqslant 0$成立.其中, $\dfrac{p_{n}+xp_{n-1}}{q_{n}+xq_{n-1}}$落在以$\dfrac{p_{n}}{q_{n}}$和$\dfrac{p_{n}+p_{n-1}}{q_{n}+q_{n-1}}$为端点的$n$阶柱集内,求和是对所有$n$阶柱集求和.\par
\par
\textbf{引理 2.  }\textsuperscript{\cite{Iosifescu}}
\begin{equation}\label{lemma2}
    \left|f_{n}^{'}(x)\right|<\frac{m}{2^{n-3}}+4M
\end{equation}
对任意$n\geqslant0$成立.\par
\par
\textbf{引理 3.  }\textsuperscript{\cite{Iosifescu}}
若
\begin{equation}\label{lemma3eq1}
    \frac{t}{1+x}<f_{n}(x)<\frac{T}{1+x},
\end{equation}
则
\begin{equation}\label{lemma3eq2}
    \frac{t}{1+x}<f_{n+1}(x)<\frac{T}{1+x},
\end{equation}
\par
\textbf{引理 4.  }\textsuperscript{\cite{Iosifescu}}
\begin{equation}\label{lemma4}
    \int\limits_{0}^{1}f_{n}(z)dz=\int\limits_{0}^{1}f_{0}(z)dz
\end{equation}
对任意$n\geqslant0$成立.\par
\par
回到\textbf{Kuzmin定理}的证明,由定理条件容易知道,存在正数$g_0,G_0,m_0=m$使得
\begin{equation}\label{Kuzmintherom1}
    \frac{g_0}{1+x}<f_0(x)<\frac{G_0}{1+x},\quad|f_0'(x)|<m_0.
\end{equation}
之后,将证明,由(\ref{Kuzmintherom1})可以推得,对充分大的$n$,存在正数$g_1,G_1,m_1$使得
$$\frac{g_1}{1+x}<f_n(x)<\frac{G_1}{1+x},\quad|f_n'(x)|<m_1,$$
并且若记$\delta=1-\dfrac{\ln2}{2}<1$,还有
\begin{align*}
    g_0<g_1<G_1<G_0, \\
    G_1-g_1<(G_0-g_0)\delta+2^{-n+2}(m_0+G_0).
\end{align*}
重复这样的步骤,可以证得,对$r=1,2\cdots$,存在正数$g_r,G_r,m_r$,使得
\begin{equation}\label{Kuzmintherom2}\begin{aligned}
         & \frac{g_r}{1+x}<f_{rn}(x)<\frac{G_r}{1+x},\quad |f_{rn}^{\prime}(x)|<n_{r}. \\
         & g_{r-1}<g_{r}<G_{r}<G_{r-1},                                                \\
         & G_r-g_r<\delta(G_{r-1}-g_{r-1})+2^{-n+2}(m_{r-1}+G_{r-1}).
    \end{aligned}\end{equation}
由\textbf{引理 3},可以取$m_r = \dfrac{m}{2^{rn-3}} + 4M$,并且对足够大的$n$,有$m_r<5M$.
代回(\ref{Kuzmintherom2})可得,存在常数$\lambda>0$,由$M,m$决定的$A_0>0$,使得
$$G_n-g_n<A_0e^{-\lambda n}.$$
记$\lim\limits_{n\to\infty}G_n=\lim\limits_{n\to\infty}g_n=a$,即有
\begin{equation}\label{Kuzmintherom3}
    \Big|f_{n^2}(x)-\frac{a}{1+x}\Big|<A_0 e^{-\lambda n},\quad \lim_{n\to\infty} \int\limits_0^1f_{n^2}(z)dz=a\ln2.
\end{equation}
由\textbf{引理 4},可知$a=\dfrac{1}{\ln2}\int_{0}^{1}f_0(z)\mathrm{d}z$.
由\textbf{引理 3},以及(\ref{Kuzmintherom3}),对足够大的$n$,都有
$$\frac{a-2A_0 e^{-\lambda\sqrt{n}}}{1+x}<f_{n}(x)<\frac{a+2A_0 e^{-\lambda\sqrt{n}}}{1+x}.$$
只需要选取足够大的常数$A\geqslant A_0$使得对有限个比较小的$n$使得要证的式子(\ref{Kuzmin})成立即可.\qed
\par
\par
下面我们运用\textbf{Kuzmin 定理}解决Gauss问题:\par
\textbf{Gauss问题的解}\textsuperscript{\cite{Kuzmin,Khinchin,Iosifescu}}\textbf{:  }
在\textbf{Kuzmin 定理}中,取$f_n(x)$为$m^{'}_n(x)$.由于$m_0(x)=x,f_0(x)=1,\forall x\in[0,1]$,故条件2、3是成立的.\par
由于$\tau^{n}(x)=\dfrac{1}{a_{n+1}+\tau^{n+1}(x)}$,从而$\tau^{n+1}(x)\leqslant x$当且仅当存在正整数$k$使得$\dfrac{1}{k+x}\leqslant\tau^{n}(x)\leqslant\dfrac{1}{k}$.从而有
$$m_{n+1}(x)=\sum_{k=1}^{\infty}m_n\left(\frac{1}{k}\right)-m_n\left(\frac{1}{k+x}\right).$$
形式上两边同时求导,可知
$$m_{n+1}^{\prime}(x)=\sum_{k=1}^{\infty}\frac{1}{(k+x)^2}m_n^{\prime}\left(\frac{1}{k+x}\right).$$
用归纳法可得$m^{'}_n(x)$一致有界,而由$\dfrac{1}{(k+x)^2}$单调趋于零,右边的求和一致收敛,上面式子是真实成立的.从而条件1也成立,从而有
$$\Big|m_n'(x)-\frac{1}{(1+x)\ln2}\Big|<\theta Ae^{-\lambda \sqrt{n}},$$
两边积分,得
$$\Big|m_n(x)-\frac{\ln{(1+x)}}{\ln{2}}\Big|\leqslant\int_0^1\Big|m_n'(x)-\frac{1}{(1+x)\ln2}\Big|\mathrm{d}x<\theta Ae^{-\lambda \sqrt{n}}.$$
这就为Gauss问题给出了一个解.


\subsection{连分数动力系统的混合性}
在这一节中,我们用\textbf{Kuzmin 定理}证明连分数动力系统的强混合性,以及连分数部分商序列的$\psi$-混合性,先介绍三个引理.\par
\textbf{引理 3.4.1.  }\textsuperscript{\cite{Liu_Peng,Mixing_lemma}}
设$E$为$k$阶柱集$\{x\in I|a_1(x)=i_1,a_2(x)=i_2,\cdots,a_k(x)=i_k\}$.则$\forall F\in\mathcal{B}_I$,有
\begin{equation}\label{mixinglemma1.1}
    \gamma(E\bigcap \tau^{-n-k}F)=\gamma(E)\gamma(F)(1+O(p^{\sqrt{n}})),\quad 0\leqslant p<1.
\end{equation}
\par
\textbf{证明:  }
记
$$F_n(u)=\{x\in I|\tau^{n+k}(x)<u\},$$
则$\tau^{-n-k}F_n(u)=[0,u).$
再记
$$f_n(u)=\dfrac{\gamma(E\cap F_n(u))}{\gamma(E)}.$$
由于$\tau^{n+k}(x)=\dfrac{1}{a_{n+k+1}+\tau^{n+k+1}(x)}$,有
\begin{equation}\label{mixinglemma1.2}
    f_{n+1}(u)=\sum_{m=1}^\infty f_{n}\left(\dfrac{1}{m}\right)-f_{n}\left(\dfrac{1}{m+u}\right).
\end{equation}\par
由于$E\cap F_0(u)=\{x\in I|a_1(x)=i_1,a_2(x)=i_2,\cdots,a_k(x)=i_k,\tau^{k}(x)<u\},$可知
$$\gamma(E\cap F_0(u))=(-1)^k\frac{1}{\ln2}\left(\ln\left(1+\frac{p_k+up_{k-1}}{q_k+uq_{k-1}}\right)-\ln\left(1+\frac{p_k}{q_k}\right)\right)$$
又
$$\gamma(E)=(-1)^k\frac{1}{\ln2}\left(\ln\left(1+\frac{p_k+p_{k-1}}{q_k+q_{k-1}}\right)-\ln\left(1+\frac{p_n}{q_n}\right)\right)$$
类似上一节中Gauss问题的处理方式,可以验证
\begin{enumerate}
    \item $\{f_n\}_{n\in\mathbb{N}}$满足\textbf{Kuzmin 定理}的条件2、3.
    \item (\ref{mixinglemma1.2})的两边形式上同时求导后,右边的求和是一致收敛的,从而上面的式子成立.故$\{f_n\}_{n\in\mathbb{N}}$满足\textbf{Kuzmin 定理}的条件 1.
\end{enumerate}\par
由\textbf{Kuzmin 定理},可知
$$|f_n'(u)-\frac{1}{(\ln 2)(1+u)}|\leqslant \theta A e^{-\lambda \sqrt{n}}.$$
上式两边同时在可测集$F$上积分,可知\\
\begin{align*}
    \left|\frac{\gamma(E\cap \tau^{-n-k}F)}{\gamma(E)}-\gamma(F)\right| & =\left|\int_F f_n(u)\mathrm{d}u-\int_F \frac{1}{(\ln 2)(1+u)}\mathrm{d}u\right| \\
                                                                        & \leqslant\int_F\left|f_n'(u)-\frac{1}{(\ln 2)(1+u)}\right|\mathrm{d}u           \\
                                                                        & \leqslant\theta A e^{-\lambda\sqrt{n}}\lambda(F)=\gamma(F)O(p^{\sqrt{n}}).
\end{align*}
上式最后一个等号用到了Gauss测度$\gamma$和Lebesgue测度$\lambda$的等价性.这也就说明了(\ref{mixinglemma1.1})成立.\qed
\par
\textbf{引理 3.4.2. }\textsuperscript{\cite{Liu_Peng,Mixing_lemma}}
记$\dfrac{p_k}{q_k}<\dfrac{\bar p_k}{\bar q_k}$是两个$k$阶收敛因子.$E=\left(\dfrac{p_k}{q_k},\dfrac{\bar p_k}{\bar q_k}\right)$,则$\forall F\in\mathcal{B}_I$,仍然成立(\ref{mixinglemma1.1}):
$$\gamma(E\bigcap \tau^{-n-k}F)=\gamma(E)\gamma(F)(1+O(p^{\sqrt{n}})),\quad 0\leqslant p<1.$$
\par\textbf{证明:  }
对区间$\left(\dfrac{p_k}{q_k},\dfrac{\bar p_k}{\bar q_k}\right)$,若不考虑在两个端点上的区别,其为至多可数个$k$阶柱集$\{x\in I|a_1(x)=i_1,a_2(x)=i_2,\cdots,a_k(x)=i_k\}$的无交并,记为
$$E=\left(\dfrac{p_k}{q_k},\dfrac{\bar p_k}{\bar q_k}\right)=\sum\limits_{i=1}^{\infty}E_i.$$
由\textbf{引理 3.4.1},以及$\sum\limits_{i=1}^{\infty}\gamma(E_i\bigcap\tau^{-n-k}F),\sum\limits_{i=1}^{\infty}\gamma(E_i)$的收敛性,有
\begin{align*}
    \gamma(\sum\limits_{i=1}^{\infty}E_i\bigcap\tau^{-n-k}F) & =\sum\limits_{i=1}^{\infty}\gamma(E_i\bigcap\tau^{-n-k}F))                                                \\
                                                             & =\sum\limits_{i=1}^{\infty}\gamma(E_i)\gamma(F)(1+O(p^{\sqrt{n}}))=\gamma(E)\gamma(F)(1+O(p^{\sqrt{n}})).
\end{align*}
这就完成了证明.\qed
\par
\textbf{引理 3.4.3.}\textsuperscript{\cite{Liu_Peng,Mixing_lemma}}
记$E$是$\mathcal{B}_I$中的任意开集,则$\forall F\in\mathcal{B}_I$,成立
\begin{equation}\label{mixinglemma3.1}
    \gamma(E\cap T^{-n}F)=\gamma(E)\gamma(F)+\gamma(F)O(p^{\sqrt{n}}),\quad 0\leqslant p<1.
\end{equation}
\par
\textbf{证明:  }
由\textbf{命题 2.1.8}, (\ref{klength})可知,任意$k$阶柱集的测度不超过$\dfrac{1}{2^{k-1}}$.\par
先考虑$E$为开区间$(a,b)$的情况,将$E=(a,b)$用三个区间覆盖.
$$E_0=\left(\dfrac{p_k}{q_k},\dfrac{\bar p_k}{\bar q_k}\right)$$
$$E_1=\left[\dfrac{p_k+p_{k-1}}{q_k+q_{k-1}},\dfrac{p_k}{q_k}\right]k\text{为奇数时,或}\left[\dfrac{p_k}{q_k},\dfrac{p_k+p_{k-1}}{q_k+q_{k-1}}\right]k\text{为偶数时};$$
$$E_2=\left[\dfrac{\bar p_k+\bar p_{k-1}}{\bar q_k+\bar q_{k-1}},\dfrac{\bar p_k}{\bar q_k}\right]k\text{为奇数时,或}\left[\dfrac{\bar p_k}{\bar q_k},\dfrac{\bar p_k+\bar p_{k-1}}{\bar q_k+\bar q_{k-1}}\right]k\text{为偶数时}.$$
其中$x\in E_1,y\in E_2$,从而$E=(a,b)\subset(E_1\bigcup E_2\bigcup E_3)$.
\begin{align*}
     & \left|\gamma(E\cap \tau^{-n}F)-\gamma(E)\gamma(F)\right|                                                                                                  \\
     & =\left|\gamma(E_0\cap \tau^{-n}F)-\gamma(E_0)\gamma(F)+\gamma(E\setminus E_0 \cap \tau^{-n}F)-\gamma(E\setminus E_0)\gamma(F)\right|                      \\
     & \leqslant\left|\gamma(E_0\cap \tau^{-n}F)-\gamma(E_0)\gamma(F)\right|+\left|\gamma(E\setminus E_0 \cap \tau^{-n}F)-\gamma(E\setminus E_0)\gamma(F)\right| \\
     & \leqslant\left|\gamma(E_0\cap \tau^{-n}F)-\gamma(E_0)\gamma(F)\right|+\gamma(F)(\gamma(E_1)+\gamma(E_2)).
\end{align*}
取$k=\lfloor \dfrac{n}{2}\rfloor$,由\textbf{引理 3.4.2}得
\begin{equation}\label{mixinglemma3.2}
    \begin{aligned}
        \left|\gamma(E_0\cap \tau^{-n}F)-\gamma(E_0)\gamma(F)\right| & \leqslant\left|\gamma(E_0\cap \tau^{-(n-k)-k}(F)-\gamma(E_0)\gamma(F)\right| \\
                                                                     & =\gamma(E_0)\gamma(F)(1+O(p^{\sqrt{n-k}}))                                   \\
                                                                     & =\gamma(E_0)\gamma(F)(1+O(p^{\sqrt{\frac{n}{2}}})).
    \end{aligned}
\end{equation}
注意到,对$0\leqslant p<1$有$0\leqslant p^{\sqrt{\frac{1}{2}}}<1$.故若记$\bar{p}=p^{\sqrt{\frac{1}{2}}}$则有, $O(p^{\sqrt{\frac{n}{2}}})=O(\bar{p}^{\sqrt{n}})$.为使记号简单,仍记(\ref{mixinglemma3.2})中的$O(p^{\sqrt{\frac{n}{2}}})$为$O(p^{\sqrt{n}})$.\par
由Gauss测度$\gamma$与Lebesgue测度$\lambda$的等价性$$\gamma(E_1)+\gamma(E_2)\leqslant\frac{1}{\ln2}(\lambda(E_1)+\lambda(E_2))=O(\frac{1}{2^{k}}).$$
并且
\begin{align*}
     & \left|\gamma(E)\gamma(F)(1+O(p^{\sqrt{n}}))-\gamma(E_0)\gamma(F)(1+O(p^{\sqrt{n}}))\right| \\
     & \leqslant(\gamma(E_1)+\gamma(E_2))\gamma(F)(1+O(p^{\sqrt{n}}))                             \\
     & =\gamma(F)O(\frac{1}{2^{k}}),
\end{align*}
从而有
$$\left|\gamma(E\cap \tau^{-n}F)-\gamma(E)\gamma(F)\right|\leqslant\gamma(E)\gamma(F)(1+O(p^{\sqrt{n}}))+\gamma(F)O(\frac{1}{2^{k}}).$$
由$k=\lfloor \dfrac{n}{2}\rfloor$,可知$O(\dfrac{1}{2^{k}})=O(p^{\sqrt{n}})$.\par
因此, (\ref{mixinglemma1.1})可写为
$$\gamma(E\cap T^{-n}F)=\gamma(E)\gamma(F)+\gamma(F)O(q^{\sqrt{n}}).$$
这就完成了$E$为开区间$(a,b)$的情况.\par
现设$E$为开集,由于其可写为可数个开区间的无交并,故也成立(\ref{mixinglemma3.1}).\qed
\par
\textbf{命题 3.4.4.}\textsuperscript{\cite{Liu_Peng,Mixing_lemma}}
连分数动力系统$(I,\mathcal{B}_I,\tau,\gamma)$是强混合的.
\par

\textbf{证明:  }
由于对任意小的$\varepsilon>0$, $\mathcal{B}_I$中的可测集$A,B$,存在开集$E,F$使得$$\lambda(A\Delta E)<\frac{1}{\ln 2}\varepsilon,\quad \lambda(B\Delta F)<\frac{1}{\ln 2}2\varepsilon.$$
此时有:
\begin{align*}
    |\gamma(A\bigcap T^{-n}B)-\gamma(E\bigcap T^{-n}B)| & <\varepsilon,  \\
    |\gamma(E\bigcap T^{-n}B)-\gamma(E\bigcap T^{-n}F)| & <\varepsilon,  \\
    |\gamma(E)\gamma(F)-\gamma(A)\gamma(B)|             & <2\varepsilon.
\end{align*}
由\textbf{引理 3.4.3},对充分大的$n$有,
$$|\gamma(E\bigcap T^{-n}F)-\gamma(E)\gamma(F)|<\varepsilon.$$
从而有
\begin{align*}
    |\gamma(A\bigcap T^{-n}B)-\gamma(A)\gamma(B)| & \leqslant|\gamma(A\bigcap T^{-n}B)-\gamma(E\bigcap T^{-n}B)| \\
                                                  & +|\gamma(E\bigcap T^{-n}A)-\gamma(E\bigcap T^{-n}F)|         \\
                                                  & +|\gamma(E\bigcap T^{-n}B)-\gamma(E)\gamma(F)|               \\
                                                  & +|\gamma(E)\gamma(F)-\gamma(A)\gamma(B)|                     \\
                                                  & <5\varepsilon
\end{align*}
这就证明了连分数动力系统的混合性.\qed
\par
根据\textbf{命题 3.4.4.}的证明,可以立即得到连分数动力系统部分商序列的$\psi$-混合性.\par
\textbf{命题 3.4.5.  }\textsuperscript{\cite{Partial_quotients,Mixing}}
连分数动力系统部分商序列是$\psi$-混合的.
\par
\textbf{证明:  }
由\textbf{命题 2.4.1},部分商序列是严平稳的,故要证明其$\psi$-混合性,只需要证明
$$\lim_{n\to\infty}\psi(X_{1}^{k},X_{k+n}^{\infty})=0.$$\par
而本节的\textbf{引理 3.4.1}已经说明, $\forall k\in\mathbb{N^+},\forall E\in X_{1}^{k},\forall F\in \mathcal{B},\gamma(E),\gamma(F)>0$,都有(\ref{mixinglemma1.1})即
$$\left|\frac{\gamma(E\bigcap \tau^{-n-k}F)}{\gamma(E)\gamma(F)}-1\right|=O(p^{\sqrt{n}}),\quad 0\leqslant p<1.$$
只需要取$F$使得$F\in X_{1}^{\infty}$,此时$\tau^{-n-k}F\in X_{k+n+1}^{\infty}$,结合$\gamma(\tau^{-n-k}F)=\gamma(F)$即可.\qed
\par
由连分数动力系统的混合性,我们可以证明连分数部分商序列的渐进独立性.\par
\textbf{命题 3.3.6.  }\textsuperscript{\cite{Liu_Peng}}
连分数部分商序列是渐进独立的.
\par
\textbf{证明:  }
对任意正整数$i,j$记
$A=\{x\in I|a_1=i\},B=\{x\in I|a_1=j\}$
由连分数动力系统的强混合性,可知
$$\lim\limits_{n\to \infty}\left|\gamma(A\cap\tau^{-n}B)-\gamma(A)\gamma(B)\right|=0,$$
则$\lim\limits_{n\to \infty}\gamma(A\cap\tau^{-n}B)=\gamma(A)\gamma(B).$
即
\begin{align*}
    \lim\limits_{n\to \infty}\gamma\{x\in I|a_1(x)=i,a_{n+1}(x)=j\} & =\gamma(\{x\in I|a_1=i\})\gamma(\{x\in I|a_1=j\})      \\
                                                                    & =\gamma(\{x\in I|a_1=i\})\gamma(\{x\in I|a_{n+1}=j\}).
\end{align*}
这就说明了连分数部分商序列的渐近独立性.\qed

\sectionbreak
% ---------------------------------------------------------------------------- %
\section{连分数动力系统的遍历性}
\subsection{遍历性简介}
\textbf{定义 4.1.1.  }
记$(X,\mathcal{B},\mu,T)$是一个保测系统.若对满足$T^{-1}A=A$的任意可测集$A\in\mathcal{B}$都有$\mu(A)=0$或$1$,则称保测变换$T:X\to X$为遍历的, $(X,\mathcal{B},\mu,T)$为遍历动力系统.
\par
遍历性有以下等价条件
\begin{enumerate}
    \item 保测变换$T$是遍历的;
    \item 若可测集$A\in\mathcal{B}$满足$\mu(T^{-1}A\Delta A)=0$,则$\mu(A)=0$或$1$;
    \item 若$A\in\mathcal{B}$且$\mu(A)>0$,则$\mu(\bigcup\limits_{i=1}^{\infty}T^{-i}A)=1$;
    \item 若$A,B\in\mathcal{B}$且$\mu(A)>0,\mu(B)>0$,则存在$n$使得$\mu(T^{-n}A\cup B)>0$.
\end{enumerate}\par
条件等价性的证明可见参考文献\cite{Ergodic_Sun,Ergodic_theory}.\par
对遍历动力系统,考虑轨道$\{T^{n}x\}_{x\in X}$,有符合直觉的下述性质.\par
\textbf{命题 4.1.1.  }\textsuperscript{\cite{Ergodic_Sun,Ergodic_theory}}
若$X$是紧致度量空间, $\mathcal{B}$为Borel  $\sigma$代数, $\mu$为Borel概率测度,且任意非空开集$A\in\mathcal{B}$都有$\mu(A)>0$,并且$T:X\to X$是连续、保测、遍历的,则几乎处处的轨道在$X$中稠密.
即
$$\mu(\{x\in X|\overline{\{T^{n}(x)|n\in\mathbb{N}\}}=X\})=1.$$
\par
\textbf{证明:  }
由$X$是紧致度量空间,可知其存在可数拓扑基,记为$\{B_n\}_{n\in\mathbb{N}^+}$.有
$$\overline{\{T^{n}(x)|n\in\mathbb{N}\}}=X\iff x\in\bigcap_{n=1}^{\infty}\bigcup_{k=0}^{\infty}T^{-k}B_n.$$
由$T^{-1}\bigcup_{k=1}^{\infty}T^{-k}B_n=\bigcup_{k=1}^{\infty}T^{-k}B_n\subset\bigcup_{k=0}^{\infty}T^{-k}B_n$,且$T$是保测的,所以
$$\mu\left(T^{-1}\bigcup_{k=1}^{\infty}T^{-k}B_n\right)=\mu\left(\bigcup_{k=1}^{\infty}T^{-k}B_n\right),$$
从而有
$$\mu\left(\left(T^{-1}\bigcup_{k=1}^{\infty}T^{-k}B_n\right)\Delta \left(\bigcup_{k=1}^{\infty}T^{-k}B_n\right)\right)=0.$$
由$T$是遍历的,由保测性的等价条件2,可知$\mu\left(\left(\bigcup_{k=1}^{\infty}T^{-k}B_n\right)\right)=0$或$1$.\par
由$T$是连续的,可知$\bigcup_{k=1}^{\infty}T^{-k}B_n$是非空开集,从而$\mu(\bigcup_{k=1}^{\infty}T^{-k}B_n)>0$,故$\mu(\bigcup_{k=1}^{\infty}T^{-k}B_n)=1$.由此可得
$$\mu\left(\bigcap_{n=1}^{\infty}\bigcup_{k=0}^{\infty}T^{-k}B_n\right)=1.$$
即几乎处处的轨道在$X$中稠密.\qed
\par
下面将介绍\textbf{Birkhoff遍历定理},它是遍历理论的基本定理.\par
\textbf{命题 4.1.2.  }\textsuperscript{\cite{Ergodic_Sun,Ergodic_theory}}\textbf{(Birkhoff遍历定理,保测情形).  }
若$(X,\mathcal{B},\mu,T)$是保测系统, $f\in L^1$,则有:
\begin{enumerate}
    \item $\dfrac{1}{n}\sum\limits_{i=0}^{n-1}f\circ T^{i}$几乎处处收敛于$L^1$函数$\bar f$.即
    $$\lim_{n\to\infty}\frac{1}{n}\sum\limits_{i=0}^{n-1}f(T^{i}(x))=\bar f(x),\quad \mathrm{a.e.}\quad x\in X.$$
    \item $\bar f(x)\circ T=\bar f(x),\quad \mathrm{a.e.}\quad x\in X.$,且有
    $$\int f\mathrm{d}m =\int \bar f\mathrm{d}m.$$
\end{enumerate}
\par
\textbf{命题 4.1.3.  }\textsuperscript{\cite{Ergodic_Sun,Ergodic_theory}}\textbf{(Birkhoff遍历定理,遍历情形).  }
若$(X,\mathcal{B},\mu,T)$是遍历系统,$f\in L^1$,则有:
\begin{equation}\label{Birkhoff1}
\lim_{n\to\infty}\frac{1}{n}\sum\limits_{i=0}^{n-1}f(T^{i}(x))=\int f\mathrm{d}m,\quad \mathrm{a.e.}\quad x\in X.
\end{equation}\par
对可测集$A\in\mathcal{B}$,由\textbf{Birkhoff遍历定理},并取$f=\mathcal{X}_A$为$A$的特征函数,则式(\ref{Birkhoff1})的左边为$x$的轨道进入$A$的频率之极限,右边为$\mu(A)$.即
\begin{equation}\label{Birkhoff2}
\mu\left(\left\{x\in X\middle|\lim_{n\to\infty}\frac{1}{n}\sum_{i=1}^{n-1}\mathcal{X}_A\circ T^{i}=\mu(A)\right\}\right)=1.
\end{equation}
\subsection{连分数动力系统的遍历性}
在本节中,我们将证明混合性与遍历性的关系,从而证明连分数动力系统的遍历性.\par
\textbf{命题 4.2.1.  }\textsuperscript{\cite{Ergodicity_and_mixing}}
强混合性蕴含弱混合性蕴含遍历性.
\par
\textbf{证明:  }
前半部分是显然的,下面证明后半部分:设$(X,\mathcal{B},\mu,T)$是弱混合的系统.若$A\in\mathcal{B}$满足$T^{-1}A=A$.则$T^{-n}(A)=A$,结合系统的弱混合性,可知,
$$\lim_{n\to\infty}{\frac{1}{n}}\sum_{i=0}^{n-1}|\mu(T^{-i}A\cap A)-\mu(A)\mu(A)|=0.$$
从而
$$\lim_{n\to\infty}{\frac{1}{n}}\sum_{i=0}^{n-1}|\mu(A)-\mu(A)\mu(A)|=0,$$
$$|\mu(A)-\mu(A)\mu(A)|=0.$$
也就是说$\mu(A)=0$或$1$.根据定义,可以知道$(X,\mathcal{B},\mu,T)$是遍历的.\qed
\par
由此,结合上一章所证的连分数动力系统的强混合性,我们可以直接得到:\par
\textbf{推论 4.2.2.  }\textsuperscript{\cite{Ergodic_theory,Ergodicity_and_mixing}}
连分数动力系统$(I,\mathcal{B}_I,\tau,\gamma)$是遍历的.
\par
根据连分数动力系统的遍历性,我们可以得到下面结论:\par
\textbf{命题 4.2.3.  }\textsuperscript{\cite{Ergodic_theory}}
$\forall i\in\mathbb{N^+}$,连分数部分商等于$i$的频率之极限在$I$上几乎处处存在且相等.
\par
\textbf{证明:  }
在(\ref{Birkhoff2})中,取$A=\{x\in I|a_1(x)=i\}=(\dfrac{1}{i+1},\dfrac{1}{i}]$,得$\gamma(A)=\dfrac{1}{\ln2}\ln\dfrac{(i+1)^2}{i(i+2)}$.
代回(\ref{Birkhoff2}),得到
$$\lim\limits_{n\to\infty}\frac{card\{k|a_k=i,1\leqslant k\leqslant n\}}{n}=\frac{1}{\ln2}\ln\dfrac{(i+1)^2}{i(i+2)},\quad\mathrm{a.e.}\quad x\in I.$$
故连分数部分商等于$i$的频率之极限在$I$上几乎处处存在且相等.\qed
\par
实际上,由此还可以得到
$$\lim\limits_{n\to\infty}\frac{card\{k|a_k\leqslant i,1\leqslant k\leqslant n\}}{n}=\sum_{j=1}^{i}\frac{1}{\ln2}\ln\dfrac{(j+1)^2}{j(j+2)},\quad\mathrm{a.e.}\quad x\in I.$$
考虑正项级数
$\sum\limits_{j=1}^{\infty}\dfrac{1}{\ln2}\ln\dfrac{(j+1)^2}{j(j+2)}=1$,从而可知连分数部分商有界的数构成一个零测集,这就是在第二章中没有给出证明的\textbf{命题2.2.5}.\par
\textbf{命题 4.2.4.  }\textsuperscript{\cite{Ergodic_theory}}
连分数部分商的几何平均数之极限在$I$上几乎处处存在且相等.
\par
\textbf{证明:  }
在\textbf{Birkhoff遍历定理}中取
$$f(x) =\ln a_1(x),$$
即有
$$\lim\limits_{n\to\infty}\sqrt[n]{(a_1(x)a_2(x)\cdots a_n(x))}=\prod\limits_{i=1}^\infty\left(\frac{(i+1)^2}{i(i+2)}\right)^{\frac{\ln{i}}{\ln{2}}},\quad\mathrm{a.e.}\quad x\in I.$$
故连分数部分商的几何平均数之极限在$I$上几乎处处存在且相等.\qed
\par
\textbf{命题 4.2.5.  }\textsuperscript{\cite{Ergodic_theory}}
连分数部分商的算术平均数之极限在$I$上几乎处处不存在.
\par
\textbf{证明:  }
注意到,如果取$f(x)=a_1(x)$,则有
$$\int\limits_{I}f\mathrm{d}x=\sum_{i=2}^{\infty}\frac{1}{i},$$发散,不满足\textbf{Birkhoff遍历定理}的条件.下面取$f_n(x)=a_1(x)\mathcal{X}_{(\frac{1}{n+1},1)},n\in\mathbb{N^+}$.\par
一方面,
$$\int\limits_{I}f_n\mathrm{d}x=\sum_{i=2}^{n+1}\frac{1}{i}.$$
另一方面,
$$f_n(T^{m}(x))=\left\{
\begin{aligned}
&a_{m+1}(x) \quad &&a_{m+1}\leqslant n\\
&0\quad &&a_{m+1}> n
\end{aligned}
\right
.$$
由\textbf{Birkhoff遍历定理}可知,$\forall n\in\mathbb{N^+}$
$$\lim_{m\to\infty}\frac{1}{m}\sum\limits_{i=0}^{n-1}f_n(T^{i}(x))=\int f_n\mathrm{d}x,\quad \mathrm{a.e.}\quad x\in I.$$
考虑$n\to\infty$的过程,即可证明
$$\lim\frac{a_1(x)+a_2(x)+\cdots+a_n(x)}{n}=\infty\quad\mathrm{a.e.}\quad x\in I.$$
故连分数部分商的算术平均数之极限在$I$上几乎处处不存在.\qed
\par
\sectionbreak

% ------------------------------------ 致谢 ------------------------------------ %
\begin{thankpage}

	原致谢已做匿名处理，感谢您阅读本论文。

\end{thankpage}

% ----------------------------------- 参考文献 ----------------------------------- %
\clearpage
\phantomsection
\addcontentsline{toc}{section}{参考文献}

\bibliographystyle{unsrt}
\bibliography{HUSTthesis}\label{bibtexref}

% ------------------------------------ 附录 ------------------------------------ %
%\clearpage
%\appendix
%\phantomsection
%\addcontentsline{toc}{section}{附录}
%\section*{附录}


\end{document}